\documentclass[11pt,reqno]{amsart}
\usepackage[T1]{fontenc}
\usepackage{lmodern}
\usepackage[letterpaper,margin=1in]{geometry}
\usepackage{microtype}
\usepackage{amsmath,amssymb,mathtools,mathrsfs}
\usepackage[hidelinks,pdfusetitle]{hyperref}
\usepackage{enumitem}
\usepackage{needspace}
\numberwithin{equation}{section}
\newtheorem{theorem}{Theorem}[section]
\newtheorem{proposition}[theorem]{Proposition}
\newtheorem{lemma}[theorem]{Lemma}
\newtheorem{corollary}[theorem]{Corollary}
\theoremstyle{definition}
\newtheorem{definition}[theorem]{Definition}
\theoremstyle{remark}
\newtheorem{remark}[theorem]{Remark}
\newcommand{\R}{\mathbb R}
\newcommand{\Z}{\mathbb Z}
\newcommand{\N}{\mathbb N}
\newcommand{\E}{\mathbb E}
\newcommand{\1}{\mathbf 1}
\newcommand{\Conf}{\operatorname{Conf}}
\newcommand{\Var}{\operatorname{Var}}
\newcommand{\PV}{\operatorname{PV}}
\newcommand{\Ent}{\operatorname{Ent}}
\newcommand{\Sine}{\operatorname{Sine}}
\newcommand{\Id}{\mathrm{Id}}
\newcommand{\Spec}{\mathsf S_P}
\newcommand{\Leb}{\operatorname{Leb}}
\DeclareMathOperator{\Law}{Law}
\newcommand{\dd}{\,d}
\newcommand{\gap}{\mathfrak g}
\newcommand{\calB}{\mathcal B}
\newcommand{\calI}{\mathcal I}
\newcommand{\calE}{\mathcal E}
\newcommand{\calF}{\mathcal F}
\newcommand{\calG}{\mathcal G}
\newcommand{\norm}[1]{\left\lVert #1\right\rVert}
\newcommand{\abs}[1]{\left\lvert #1\right\rvert}
\DeclareMathOperator{\supp}{supp}
\DeclareMathOperator{\diver}{div}
\allowdisplaybreaks[2]
\setlength{\emergencystretch}{1.5em}
\setlist[enumerate]{leftmargin=*,itemsep=3pt,topsep=5pt}
\title[Uniqueness of the logarithmic DLR state]{Uniqueness of the stationary logarithmic DLR state}
\date{Research manuscript draft\\9 September 2026; exposition revised 10 September 2026}
\subjclass[2020]{60G55, 60K35, 82B21, 60B20}
\keywords{Logarithmic gas, DLR equations, Sine beta process, hyperuniformity, Palm measure, entropy}

\begin{document}
\begin{abstract}
Fix an inverse temperature $\beta>0$. For a stationary simple point
process $P$ on $\R$ with intensity one, write
$v_P(t)=\Var_P N_{(0,t]}$. We prove that the canonical logarithmic
Dobrushin--Lanford--Ruelle equations have at most one such solution
satisfying
\[
v_P(t)\le C_Pt\quad(t>0),\qquad
\lim_{t\to\infty}\frac{v_P(t)}t=0,\qquad
\int_1^\infty\frac{v_P(t)}{t^2}\dd t<\infty.
\]
The proof compares the conditional distributions of two candidate
processes with the same finite-dimensional Gibbs measure: a logarithmic
gas with lattice exterior particles and a quadratic potential of
vanishing strength. A concentration inequality for averages of
consecutive gaps controls the difference of expectations by relative
entropy per particle. The integrated variance bound controls the
exterior-force contribution in probability, and the sublinear variance
bound makes the quadratic displacement contribution tend to zero.
We then derive the variance assumptions from finite renormalized
logarithmic electric energy, using a stationary joint distribution of
the configuration and an electric field and a covariance spectral
estimate. With the energy normalization specified below, this identifies
every stationary finite-energy canonical DLR solution with $\Sine_\beta$.
\end{abstract}
\maketitle

\section{Introduction and statements}

The process $\Sine_\beta$ describes the microscopic bulk limit of
one-dimensional logarithmic gases. We use its unit-intensity
normalization throughout. Dereudre, Hardy, Lebl\'e and
Ma\"ida~\cite{DHLM} show that it satisfies canonical Dobrushin--Lanford--Ruelle
equations and that every stationary finite-energy solution of those
equations is number-rigid. Their Section~1.3 asks whether the stationary
solution is unique. The conditional universality theorem of Kuijlaars and
Mi\~na-D\'iaz~\cite{KMD} supplies a deterministic route at $\beta=2$.
For general $\beta$, Bourgade, Erd\H{o}s and Yau~\cite{BEY} develop a
comparison method for local logarithmic Gibbs measures, together with
strong particle-location estimates for convex analytic potentials.

The organizing idea of the argument below is to compare both candidate
infinite-volume laws with one common Gibbs measure with lattice exterior and then
subtract the two comparison estimates. No limit of the reference measure is needed. A static
convexity estimate converts relative entropy into control of averaged
gap statistics. The integrated variance bound gives tightness of the conditional squared
exterior-force difference. Sublinear variance controls the quadratic
displacement introduced by the additional potential. The comparison theorem is stated directly in terms of
these count-variance hypotheses and is independent of the electric-field
argument used later to verify them.

Assiotis~\cite[Theorem~1.1]{Assiotis} proves stationary uniqueness for
every $\beta>0$ in the same finite-energy class, without an ergodicity
assumption. His process energy is exactly $(2\pi)^{-1}\E_PW$;
Lemma~\ref{lem:energy-input} checks the definitions, and
Remark~\ref{rem:dlr-conventions} compares the DLR formulations.
His proof combines relative-entropy comparison with a circular
$\beta$ ensemble, cyclic-gap transport, the circular Palm limit,
and Palm inversion. The argument below shares the entropy and Palm
framework and provides an alternative comparison proof of that
finite-energy uniqueness result. Its reference is a Gibbs measure with lattice exterior
and a weak quadratic potential, which cancels between the two
candidates without identifying its limit. Theorem~\ref{thm:comparison}
isolates the discrepancy assumptions sufficient for this comparison.

\subsection{Configurations, notation, and the DLR specification}
\label{subsec:notation}

Let $\Conf(\R)$ denote the space of locally finite simple counting
measures on $\R$, with the Borel sigma-algebra of the vague topology.
For $\gamma\in\Conf(\R)$ and a Borel set $A$, write
\[
N_A(\gamma)=\gamma(A),\qquad
\gamma_A=\gamma|_A,\qquad
(\theta_a\gamma)(A)=\gamma(A+a)\quad(a\in\R).
\]
Thus $\theta_a$ moves a point at $a$ to the origin. A probability measure
$P$ on $\Conf(\R)$ is stationary if $(\theta_a)_\#P=P$ for every $a$.
Here $T_\#\nu(A)=\nu(T^{-1}(A))$ denotes the pushforward of a
measure $\nu$ under a measurable map $T$. For a probability measure,
it is the distribution of $T(X)$ when $X$ has law $\nu$.
We also write $\Law_\nu(X)$ for the distribution of a random variable
$X$ under $\nu$.
It has intensity one if
\begin{equation}\label{eq:intensity-one}
\E_PN_A=\Leb(A)
\qquad\text{for every bounded Borel set }A\subset\R.
\end{equation}
Here $\Leb$ is Lebesgue measure, and $\delta_x$ denotes the unit
point mass at $x$. We use $\N=\{1,2,\ldots\}$. We omit the argument $\gamma$ from
counts when the configuration is clear. For a probability measure $\nu$,
we use $\nu(f)=\int f\dd\nu$ and $\Var_\nu f=\nu(f^2)-\nu(f)^2$.
For a stationary intensity-one process, set
\begin{equation}\label{eq:variance-definition}
v_P(t)=\Var_PN_{(0,t]}
      =\E_P\bigl[(N_{(0,t]}-t)^2\bigr],\qquad t>0.
\end{equation}
For vectors, $|\cdot|$ is the Euclidean norm; $\Id_K$ is the
$K\times K$ identity matrix. A generic constant $C$ may change from
one occurrence to the next. Constants carrying a subscript record
additional dependencies. In estimates for $K\to\infty$, constants are
independent of $K$ and of the confinement parameter $T$ unless stated
otherwise; dependence on fixed $\beta,r,G$, or on a fixed candidate law,
is allowed. We write $X_K=o_P(a_K)$ for $X_K/a_K\to0$ in
$P$-probability, and $x_+=\max\{x,0\}$. The notation
$a_K\asymp b_K$ means $c_1b_K\le a_K\le c_2b_K$ for positive
constants $c_1,c_2$ independent of $K$.

We specify the conditional distribution meant by the canonical DLR
equations. Let $I=(a,b)$ be a bounded interval, let $n\ge1$, and define
\[
\Omega_{I,n}=\{x\in\R^n:a<x_1<\cdots<x_n<b\}.
\]
For an exterior configuration $\xi$ supported in $I^c$, choose a
reference vector $y\in\Omega_{I,n}$ and define the relative Hamiltonian
\begin{equation}\label{eq:dlr-relative-energy}
\begin{split}
\mathscr H_I^\xi(x;y)
={}&-\sum_{1\le i<j\le n}
       \log\frac{|x_j-x_i|}{|y_j-y_i|}\\
&-\lim_{R\to\infty}
  \sum_{\substack{p\in\xi\\|p|\le R}}\sum_{i=1}^n
       \log\frac{|x_i-p|}{|y_i-p|}.
\end{split}
\end{equation}
The cutoff in this definition is symmetric about the spatial origin.
Whenever a principal-value sum $\PV\sum_p$ occurs below, it means this
limit with $|p|\le R$ in the coordinates then in use; any explicitly
excluded particles are omitted before taking the limit. The relative
energies and the corresponding normalized kernel are required to exist
for almost every exterior and particle number arising under the process.
The kernel is
\begin{equation}\label{eq:dlr-kernel}
\begin{aligned}
\kappa_{I,n}^\beta(\xi;\dd x)
 &=\frac{e^{-\beta\mathscr H_I^\xi(x;y)}}
         {Z_{I,n}^\beta(\xi;y)}\,\1_{\Omega_{I,n}}(x)\dd x,\\
Z_{I,n}^\beta(\xi;y)
 &=\int_{\Omega_{I,n}}e^{-\beta\mathscr H_I^\xi(x;y)}\dd x
 \in(0,\infty).
\end{aligned}
\end{equation}
Changing $y$ adds a constant independent of $x$ to the Hamiltonian and
therefore does not change the kernel. For $n=0$, the kernel is the point
mass on the empty configuration.

Writing $\mathscr G_I=\sigma(\gamma_{I^c},N_I)$, the canonical DLR
property at inverse temperature $\beta$ is the identity
\begin{equation}\label{eq:dlr-identity}
\E_P[f(\gamma)\mid\mathscr G_I]
 =\int f\!\left(\gamma_{I^c}+\sum_{i=1}^{N_I}\delta_{x_i}\right)
       \kappa_{I,N_I}^\beta(\gamma_{I^c};\dd x)
 \quad P\text{-almost surely}
\end{equation}
for every bounded measurable $f:\Conf(\R)\to\R$ and every bounded
interval $I$. This is the ordered-coordinate form of the relative-energy
specification in \cite[Section~2.1]{DHLM}.

\Needspace{10\baselineskip}
\begin{samepage}
\begin{theorem}[Uniqueness under discrepancy control]\label{thm:comparison}
Fix $\beta>0$. There is at most one stationary simple point process of
intensity one satisfying the canonical logarithmic DLR equations and the
following three conditions:
\begin{align}
v_P(t)&\le C_Pt \quad(t>0),\label{eq:linear}\\
v_P(t)&=o(t)\quad(t\to\infty),\label{eq:hyper}\\
\int_1^\infty v_P(t)t^{-2}\dd t&<\infty.\label{eq:integrated}
\end{align}
The constant $C_P$ need not be uniform over candidate processes.
\end{theorem}
\end{samepage}

For clarity, all three conditions are retained in the statement.
Section~\ref{sec:energy} proves them together for each ergodic finite-energy
DLR solution. The comparison proof itself requires no ergodicity.

\begin{samepage}
\begin{theorem}[Finite-energy uniqueness]\label{thm:finite-energy}
Fix $\beta>0$, and let $W$ be the renormalized logarithmic energy
defined in \eqref{eq:renormalized-energy}, with the inner truncation and
unit background specified in Section~\ref{sec:energy}.
If $P$ is stationary, satisfies the canonical logarithmic DLR
property \eqref{eq:dlr-identity}, and $\E_P W<\infty$, then
\[
P=\Sine_\beta.
\]
\end{theorem}
\end{samepage}

Theorem~\ref{thm:finite-energy} follows from
Theorem~\ref{thm:comparison}, Proposition~\ref{prop:energy-discrepancy},
the ergodic DLR decomposition in \cite[Proposition~3.13]{DHLM}, and the
finite-energy existence input of \cite[Corollary~1.2]{LS}.
The energy condition uses a background of intensity one. Mean intensity
one by itself is insufficient for uniqueness: the equal mixture of the
vacuum and the image of $\Sine_\beta$ under $x\mapsto x/2$ is a
stationary canonical DLR solution of intensity one and infinite energy
\cite[Proposition~9.4]{Assiotis}. Indeed, if $P_{\rm mix}$ denotes that mixture and $\Lambda$ its
component intensity, then
\[
P(\Lambda=0)=P(\Lambda=2)=\tfrac12,\qquad
\Var_{P_{\rm mix}}N_{(0,t]}
\ge\Var(\Lambda t)=t^2.
\]
It therefore also fails \eqref{eq:linear}.

\subsection{Organization and main estimates}

Section~\ref{sec:entropy} defines the finite-dimensional reference measure,
the gap average, relative entropy, and relative Fisher information before
stating the comparison inequalities \eqref{eq:entropy-gap} and
\eqref{eq:fisher-gap}. The reference has lattice exterior particles and
a quadratic potential of strength $T^{-1}$. Section~\ref{sec:palm} collects the
counting and Palm facts needed for random particle-ended intervals.
Sections~\ref{sec:blocks} and \ref{sec:forces} control, respectively, the
confinement cost and the exterior-force contribution. These estimates
are assembled in Section~\ref{sec:proof-main}. The electric-field
argument in Section~\ref{sec:energy} verifies the hypotheses of
Theorem~\ref{thm:comparison} for finite-energy states.


The three variance hypotheses have distinct roles in the comparison:
\begin{center}
\renewcommand{\arraystretch}{1.2}
\begin{tabular}{@{}p{0.32\textwidth}p{0.62\textwidth}@{}}
\hline
Assumption & Main role \\
\hline
Linear bound \eqref{eq:linear}
& Counting regularity and integrable Palm count envelopes
(Section~\ref{sec:palm}). \\
Sublinear growth \eqref{eq:hyper}
& Endpoint-length control and negligible quadratic displacement
(Section~\ref{sec:blocks}). \\
Integrated variance \eqref{eq:integrated}
& Half-line discrepancy and exterior-force envelopes
(Sections~\ref{sec:palm} and \ref{sec:forces}). \\
\hline
\end{tabular}
\end{center}
Proposition~\ref{prop:energy-discrepancy} collects the three conclusions
for ergodic finite-energy components. Appendix~\ref{app:convex} supplies
the convex-domain and weak Sobolev justification used in the static
comparison estimate.

\section{A static comparison inequality for gap statistics}
\label{sec:entropy}

Fix $K\ge1$, write $L=K+1$, and let
\[
\Omega_K=\{u\in\R^K:0<u_1<\cdots<u_K<L\}.
\]
Set $u_0=0$, $u_{K+1}=L$, and $g_i(u)=u_{i+1}-u_i$ for
$0\le i\le K$. The reference exterior is
\[
\xi_K^{\rm lat}=\{j\in\Z:j\le0\text{ or }j\ge L\}.
\]
With any fixed $r_*\in(0,L)$, define
\begin{equation}\label{eq:lattice-energy}
H_K^{\rm lat}(u)=
-\sum_{1\le i<j\le K}\log(u_j-u_i)
-\sum_{i=1}^K\PV\sum_{p\in\xi_K^{\rm lat}}
       \log\abs{\frac{u_i-p}{r_*-p}}.
\end{equation}
Changing $r_*$ changes only an additive constant. Euler's product for
the gamma function gives the following exact finite expression:
\begin{equation}\label{eq:lattice-gamma-energy}
H_K^{\rm lat}(u)=
-\sum_{1\le i<j\le K}\log(u_j-u_i)
+\sum_{i=1}^K\log
 \frac{\Gamma(u_i)\Gamma(L-u_i)}{\Gamma(r_*)\Gamma(L-r_*)}.
\end{equation}
Here $\Gamma(z)=\int_0^\infty t^{z-1}e^{-t}\dd t$ for $z>0$.
The paired exterior series converges, and its derivatives of order at
least two converge locally uniformly. The endpoint singularities are
$-\log u_i-\log(L-u_i)$.
For $T>0$ put
\begin{equation}\label{eq:reference}
\omega_{K,T}(\dd u)=Z_{K,T}^{-1}e^{-\Phi_{K,T}(u)}\dd u,\qquad
\Phi_{K,T}(u)=\beta H_K^{\rm lat}(u)
                  +\frac1{2T}\sum_{i=1}^K(u_i-i)^2.
\end{equation}
Here and below $\dd u$ denotes $K$-dimensional Lebesgue measure, and
\begin{equation}\label{eq:reference-normalizer}
Z_{K,T}=\int_{\Omega_K}e^{-\Phi_{K,T}(u)}\dd u\in(0,\infty).
\end{equation}

Fix an integer $r\ge1$ and a function $G\in C_c^2((0,\infty)^r)$.
For $K>r$, define
\begin{equation}\label{eq:gap-observable}
F_K(u)=\frac1{K-r}
\sum_{i=1}^{K-r}G(g_i(u),\ldots,g_{i+r-1}(u)).
\end{equation}
This uses only gaps between interior particles. Adding or deleting at
most $m$ edge windows, with $m$ fixed, changes the normalized average
by at most $2m\norm G_\infty/(K-r-m)$ for $K>r+m$; this bound tends
to zero as $K\to\infty$.

For probability measures $\mu\ll\omega$ on $\Omega_K$, write
$h=\dd\mu/\dd\omega$. The relative entropy and relative Fisher
information are, respectively,
\begin{equation}\label{eq:entropy-fisher-definitions}
S(\mu\mid\omega)=\int h\log h\dd\omega,\qquad
I(\mu\mid\omega)=4\int\abs{\nabla\sqrt h}^2\dd\omega,
\end{equation}
with $0\log0=0$. If $\mu\not\ll\omega$, both quantities are defined
to be $+\infty$.
Here $I=+\infty$ unless $\sqrt h$ belongs to the weighted weak
Sobolev space specified in Appendix~\ref{app:convex}. For positive
locally differentiable densities this is
$\int|\nabla\log h|^2\dd\mu$.

\begin{samepage}
\begin{proposition}[Entropy and Fisher-information bounds for gap averages]\label{prop:gap-entropy}
Fix $\beta>0$, an integer $r\ge1$, and $G\in C_c^2((0,\infty)^r)$.
There is a constant $C_G<\infty$ such that, for every integer
$K\ge2r+1$, every $T>0$, and every probability measure $\mu$ on
$\Omega_K$,
\begin{align}
|\mu(F_K)-\omega_{K,T}(F_K)|^2
&\le \frac{C_G}{K}S(\mu\mid\omega_{K,T}),\label{eq:entropy-gap}\\
|\mu(F_K)-\omega_{K,T}(F_K)|^2
&\le \frac{C_GT}{K}I(\mu\mid\omega_{K,T}).\label{eq:fisher-gap}
\end{align}
More precisely, choose $B<\infty$ so that
$\supp G\subset(0,B)^r$, and put
\[
M_0=\norm G_\infty,\qquad
M_1=\max_a\norm{\partial_aG}_\infty,\qquad
M_2=\max_{a,b}\norm{\partial_{ab}G}_\infty.
\]
The constant may be chosen to depend only on
$\beta,r,B,M_0,M_1,M_2$.
\end{proposition}
\end{samepage}

\begin{proof}
Write $\Phi=\Phi_{K,T}$ and $n=K-r\ge K/2$. For $v\in\R^K$,
extend $v$ by $v_0=v_{K+1}=0$ and put $d_i=v_{i+1}-v_i$.
Nearest-neighbor interactions and the two endpoint charges give
\begin{equation}\label{eq:hessian}
v^\top\Phi''(u)v
\ge\beta\sum_{i=0}^K\frac{d_i^2}{g_i(u)^2}
   +\frac1T|v|^2.
\end{equation}
Every remaining logarithmic interaction has nonnegative Hessian.
We use the stronger gap-dependent part of this lower bound to obtain
the factor $K^{-1}$; the uniform part alone would not give it.

\smallskip\noindent\textit{First gap derivatives.}
Regard $F_K$ as a function of the gaps and set
$q_j=\partial F_K/\partial g_j$. Each gap belongs to at most $r$
windows, so
\[
q_j=\frac1n\sum_{\substack{1\le i\le n\\i\le j\le i+r-1}}
  \partial_{j-i+1}G(g_i,\ldots,g_{i+r-1}),\qquad
|q_j|\le\frac{rM_1}{n}\le\frac{2rM_1}{K}.
\]
The endpoint coefficients $q_0,q_K$ vanish. Since the support of
each derivative is contained in $\supp G$, a nonzero $q_j$ requires
$g_j\le B$. Consequently
\[
\sum_{j=0}^K q_j^2g_j^2\le\frac{4r^2M_1^2B^2}{K}.
\]
Weighted Cauchy--Schwarz and \eqref{eq:hessian} now give
\[
|\nabla F_K\cdot v|^2
=\abs{\sum_jq_jd_j}^2
\le\left(\beta^{-1}\sum_jq_j^2g_j^2\right)
     \left(\beta\sum_j\frac{d_j^2}{g_j^2}\right)
\le\frac{A_1}{K}v^\top\Phi''v,
\qquad A_1=\frac{4r^2B^2M_1^2}{\beta}.
\]
Taking the supremum over $v$ in this quadratic-form duality yields
\begin{equation}\label{eq:gradient-metric}
\nabla F_K^\top(\Phi'')^{-1}\nabla F_K\le A_1/K.
\end{equation}

\smallskip\noindent\textit{Second gap derivatives and window overlap.}
Let $w_i=(g_i,\ldots,g_{i+r-1})$ and let $J_i$ be the indicator
that $w_i$ belongs to the support of $D^2G$. Then
\[
\begin{split}
|v^\top F_K''v|
&\le\frac{M_2}{n}\sum_{i=1}^n
  J_i\left(\sum_{a=0}^{r-1}|d_{i+a}|\right)^2\\
&\le\frac{rM_2}{n}\sum_{i=1}^nJ_i
  \sum_{a=0}^{r-1}d_{i+a}^2\\
&\le\frac{rM_2B^2}{n}\sum_{i=1}^n
  \sum_{a=0}^{r-1}\frac{d_{i+a}^2}{g_{i+a}^2}
\le\frac{r^2M_2B^2}{n}\sum_{j=0}^K\frac{d_j^2}{g_j^2}.
\end{split}
\]
In the third line, every gap in an active window is at most $B$;
in the last line, every gap is counted at most $r$ times. Thus
\begin{equation}\label{eq:tilt-hessian}
|v^\top F_K''v|
\le\frac{A_2}{K}
 v^\top(\beta H_K^{{\rm lat}\,\prime\prime})v,
\qquad A_2=\frac{2r^2B^2M_2}{\beta}.
\end{equation}

\smallskip\noindent\textit{Uniformity over the tilted laws.}
For $s\in\R$, define the exponentially tilted probability measure
\[
\omega_s(\dd u)=
\frac{e^{sF_K(u)}}{\omega_{K,T}(e^{sF_K})}\,\omega_{K,T}(\dd u),
\]
Set $c_G=[2(1+A_2)]^{-1}$. The same pointwise estimate
\eqref{eq:tilt-hessian}, with no dependence on $s,K,T$, gives
for every $|s|\le c_GK$
\[
(\Phi-sF_K)''
\ge\tfrac12\beta H_K^{{\rm lat}\,\prime\prime}+T^{-1}\Id_K
\ge\tfrac12\Phi''.
\]
The Brascamp--Lieb inequality in Lemma~\ref{lem:convex-approx},
followed by \eqref{eq:gradient-metric}, therefore gives
\[
\Var_{\omega_s}F_K
\le\E_{\omega_s}\bigl[
\nabla F_K^\top((\Phi-sF_K)'')^{-1}\nabla F_K\bigr]
\le\frac{2A_1}{K}
\qquad(|s|\le c_GK).
\]
This estimate first holds on each gap-truncated simplex and passes
to $\omega_s$ with the same constants, as proved in the appendix.
If $\psi(s)=\log\omega_{K,T}(e^{sF_K})$, then
$\psi''(s)=\Var_{\omega_s}F_K$. Integrating twice from zero yields
\begin{equation}\label{eq:mgf}
\log\E_{\omega_{K,T}}
\exp\{s(F_K-\omega_{K,T}(F_K))\}
\le A_1s^2/K,\qquad |s|\le c_GK.
\end{equation}

\smallskip\noindent\textit{Entropy and Fisher information.}
The entropy inequality, used with either sign, gives
\[
|\mu(F_K)-\omega_{K,T}(F_K)|
\le\frac{S(\mu\mid\omega_{K,T})}{s}+\frac{D s}{K},
\qquad 0<s\le c_GK,\qquad D=1+A_1.
\]
Abbreviate $S=S(\mu\mid\omega_{K,T})$ for the rest of this paragraph.
When $0<S\le Dc_G^2K$, choose $s=\sqrt{SK/D}$ to obtain the
square bound $4DS/K$; the case $S=0$ follows by $s\downarrow0$.
When $S>Dc_G^2K$, use
$|\mu(F_K)-\omega_{K,T}(F_K)|\le2M_0$ instead.
Thus \eqref{eq:entropy-gap} holds, for example, with any
\[
C_G\ge\max\{4D,\,4M_0^2/(Dc_G^2)\}.
\]
When $S=+\infty$, \eqref{eq:entropy-gap} is immediate.
Finally, the uniform convexity in \eqref{eq:hessian} and the
logarithmic Sobolev part of Lemma~\ref{lem:convex-approx} imply
\begin{equation}\label{eq:lsi}
S(\mu\mid\omega_{K,T})\le\frac T2 I(\mu\mid\omega_{K,T}),
\end{equation}
which proves \eqref{eq:fisher-gap} after enlarging the constant if
necessary. The static convexity inequalities are also used in
\cite[Sections~3 and 5]{BEY}. The appendix makes their use on the
singular simplex, including the weak Fisher-information domain,
explicit for every $\beta>0$.
\end{proof}


\section{Discrepancy and Palm preliminaries}\label{sec:palm}

Throughout Sections~\ref{sec:palm}--\ref{sec:proof-main}, $P$ satisfies
Theorem~\ref{thm:comparison}. Its Palm distribution $P^0$ is the
probability measure on configurations with a point at zero characterized
by Campbell's identity
\begin{equation}\label{eq:campbell}
\E_P\sum_{p\in\gamma} f(p,\theta_p\gamma)
 =\int_\R\E_{P^0}f(x,\gamma)\dd x
\end{equation}
for every nonnegative measurable $f$ on $\R\times\Conf(\R)$.
For a rooted configuration with infinitely many points in each
direction, enumerate its points and gaps as
\begin{equation}\label{eq:palm-points-gaps}
\cdots<p_{-1}<p_0=0<p_1<\cdots,\qquad
 g_j=p_{j+1}-p_j\quad(j\in\Z).
\end{equation}
The successor shift is $T_*\gamma=\theta_{p_1}\gamma$; thus
$g_j(T_*\gamma)=g_{j+1}(\gamma)$. Palm measure is invariant under this
map: $(T_*)_\#P^0=P^0$. The almost-sure existence of the two-sided
enumeration follows from the next lemma.

\begin{samepage}
\begin{lemma}[Almost-sure exterior regularity]\label{lem:as-regularity}
There is a translation-invariant measurable set $\mathcal R$ with
$P(\mathcal R)=P^0(\mathcal R)=1$ such that, for every
$\gamma\in\mathcal R$ and every $a\in\R$,
\begin{equation}\label{eq:as-density}
\lim_{t\to\infty}\frac{N_{(a,a+t]}(\gamma)}t
 =\lim_{t\to\infty}\frac{N_{[a-t,a)}(\gamma)}t=1,
\end{equation}
and
\begin{equation}\label{eq:as-integral}
\int_1^\infty\frac{
  (N_{(a,a+t]}-t)^2+(N_{[a-t,a)}-t)^2}{t^2}\dd t<\infty.
\end{equation}
Moreover,
\begin{equation}\label{eq:reciprocal-pv}
\lim_{R\to\infty}
\sum_{\substack{p\in\gamma\\0<|p-a|\le R}}\frac1{p-a}
\quad\text{exists as a finite real number for every }a\in\R.
\end{equation}
A point at $a$, when present, is omitted from this sum.
\end{lemma}
\end{samepage}

\begin{proof}
By \eqref{eq:linear}, Chebyshev's inequality along $t=n^2$ and
Borel--Cantelli give $N_{(0,t]}/t\to1$; monotonic interpolation removes
the subsequence. The backward statement is identical.
Tonelli and \eqref{eq:integrated} give \eqref{eq:as-integral} at $a=0$.
Define the oriented cumulative discrepancy by
\begin{equation}\label{eq:oriented-discrepancy}
D_0(x)=
\begin{cases}
N_{(0,x]}-x,&x\ge0,\\
-N_{(x,0]}-x,&x<0.
\end{cases}
\end{equation}
Then $N_{(a,a+t]}-t=D_0(a+t)-D_0(a)$ for every $t>0$.
The backward discrepancy with the convention $[a-t,a)$ differs from
$D_0(a)-D_0(a-t)$ by at most two endpoint indicators. A change of
variable and comparability of $t$ and $a+t$ for large $t$ prove the
integrability at every finite $a$, simultaneously on the same event.

For the principal value, put
$D_{\rm sym}(t)=N_{(0,t]}-N_{[-t,0)}$. Then
\[
\sum_{1<|p|\le R,\ p\in\gamma}\frac1p
=\frac{D_{\rm sym}(R)}R-D_{\rm sym}(1)
  +\int_1^R\frac{D_{\rm sym}(t)}{t^2}\dd t.
\]
The first term vanishes by density one. The integral converges
absolutely by Cauchy--Schwarz and \eqref{eq:as-integral}. Add the finite
sum over $0<|p|\le1$. The same argument at any $a$ gives
\eqref{eq:reciprocal-pv}. All the resulting properties hold simultaneously
at every $a$, so their common event is translation invariant.
Equation~\eqref{eq:campbell} transfers this event to Palm measure.
\end{proof}

\begin{samepage}
\begin{lemma}[Integrable bounds for Palm counting functions]\label{lem:palm-envelope}
Under $P^0$ define
\[
C^+=\sup_{t\ge0}\frac{N_{(0,t]}}{1+t},\qquad
C^-=\sup_{t\ge0}\frac{N_{[-t,0)}}{1+t}.
\]
Then
\begin{equation}\label{eq:palm-envelope-integrability}
\E_{P^0}C^++\E_{P^0}C^-<\infty.
\end{equation}
\end{lemma}
\end{samepage}

\begin{proof}
For a spatially rooted sample put
$C_0=\sup_{t\ge0}N_{(0,t]}/(1+t)$.
A dyadic bound gives
\[
C_0\le N_{(0,1]}+2+
 \sum_{j\ge0}2^{-j}
 \abs{N_{(0,2^{j+1}]}-2^{j+1}}.
\]
Minkowski's inequality and \eqref{eq:linear} show
$\E_PC_0^2<\infty$. For every $p\in\gamma\cap(0,1]$,
$C^+(\theta_p\gamma)\le2C_0(\gamma)$.
Campbell's formula at intensity one therefore gives
\[
\E_{P^0}C^+
=\E_P\sum_{p\in\gamma\cap(0,1]}C^+(\theta_p\gamma)
\le2\E_P[N_{(0,1]}C_0]<\infty.
\]
Reflect the argument for $C^-$.
\end{proof}

\subsection{The first positive point and the preceding gap}

Use the Palm points and gaps in \eqref{eq:palm-points-gaps}.
Campbell's identity gives the following size-bias formula. If $a$ is the first positive point of a spatial sample, the law
of $\theta_a\gamma$ is
\begin{equation}\label{eq:first-positive}
\widetilde P^0(\dd\gamma)=g_{-1}(\gamma)P^0(\dd\gamma).
\end{equation}
Indeed, each gap ending at a point $p$ contributes the interval of
possible spatial origins of length $g_{-1}(\theta_p\gamma)$.
In particular $\E_{P^0}g_{-1}=1$.

\begin{samepage}
\begin{lemma}[Uniform domination of endpoint environments]
\label{lem:endpoint-domination}
Let $a_j$ be the $(j+1)$-st positive point of a spatial sample, and let
$\gap_*$ be the length of the whole gap containing its spatial origin.
For every $M>0$, every integer $j\ge0$, and every nonnegative measurable
function $\Phi$ on rooted configurations,
\begin{equation}\label{eq:endpoint-domination}
\E_P\bigl[\1_{\{\gap_*\le M\}}\Phi(\theta_{a_j}\gamma)\bigr]
\le M\E_{P^0}\Phi.
\end{equation}
Thus the restricted endpoint laws are dominated by $MP^0$, uniformly
in the number $j$ of successor shifts.
\end{lemma}
\end{samepage}

\begin{proof}
The event $\{\gap_*\le M\}$ becomes $\{g_{-1}\le M\}$ when the sample
is rooted at its first positive point. Since
$\theta_{a_j}\gamma=T_*^j(\theta_{a_0}\gamma)$,
\eqref{eq:first-positive} gives
\[
\E_P\bigl[\1_{\{\gap_*\le M\}}\Phi(\theta_{a_j}\gamma)\bigr]
=\E_{P^0}\bigl[g_{-1}\1_{\{g_{-1}\le M\}}
                    \Phi(T_*^j\gamma)\bigr]
\le M\E_{P^0}[\Phi\circ T_*^j]
=M\E_{P^0}\Phi.
\]
The last equality uses $(T_*)_\#P^0=P^0$. No invariance of the
truncating event, and no independence of that event from the shifted
environment, is required.
\end{proof}

The next lemma addresses the bias in \eqref{eq:first-positive} without
assuming ergodicity. The first positive point has the preceding-gap
weighted law, so a pointwise ergodic limit of a block average need not
itself be constant. What removes the weight is that the conditional
mean of the preceding gap on every point-shift invariant component is
one.

\begin{samepage}
\begin{lemma}[Spatial block averages identify Palm gap laws]
\label{lem:identify}
Let $r\ge1$ be an integer and let $G:(0,\infty)^r\to\R$ be bounded
and measurable. Let $a_j$ be the $(j+1)$-st positive point of a
configuration sampled under $P$, for $j\ge0$. For $K>r$, set
\begin{equation}\label{eq:spatial-gap-average}
A_{K,r}(\gamma)=\frac1{K-r}\sum_{i=1}^{K-r}
G(a_{i+1}-a_i,\ldots,a_{i+r}-a_{i+r-1}).
\end{equation}
Then
\begin{equation}\label{eq:spatial-palm-limit}
\lim_{K\to\infty}\E_PA_{K,r}
 =\E_{P^0}G(g_0,\ldots,g_{r-1}).
\end{equation}
\end{lemma}
\end{samepage}

\begin{proof}
Put $\Psi(\gamma)=G(g_0,\ldots,g_{r-1})$ and let $\calI$ be the
invariant sigma-algebra of $T_*$. Define the rooted average
\[
A_K^0=\frac1{K-r}\sum_{i=1}^{K-r}\Psi\circ T_*^i,
\qquad A_{K,r}(\gamma)=A_K^0(\theta_{a_0}\gamma).
\]
Birkhoff's theorem gives
$A_K^0\to\E_{P^0}[\Psi\mid\calI]$ $P^0$-almost surely. Shifting the
starting index by a fixed number changes a bounded average by $O(K^{-1})$.
Density one and $\E_{P^0}g_{-1}<\infty$ give
\[
\E_{P^0}[g_{-1}\mid\calI]
=\lim_{n\to\infty}\frac1n\sum_{j=0}^{n-1}g_{-1}\circ T_*^j
=\lim_{n\to\infty}\frac{p_{n-1}-p_{-1}}n=1.
\]
Here the middle equality telescopes the gaps and the last uses the
almost-sure particle density from Lemma~\ref{lem:as-regularity}.
By \eqref{eq:first-positive}, the expectation of the spatially rooted
block average is $\E_{P^0}[g_{-1}A_K^0]$. Since
$|A_K^0|\le\norm G_\infty$, dominated convergence with the integrable
weight $g_{-1}$, followed by conditioning on $\calI$, yields
\[
\lim_{K\to\infty}\E_{P^0}[g_{-1}A_K^0]
=\E_{P^0}\!\left[g_{-1}\E_{P^0}[\Psi\mid\calI]\right]
=\E_{P^0}\!\left[\E_{P^0}[g_{-1}\mid\calI]
                         \E_{P^0}[\Psi\mid\calI]\right]
=\E_{P^0}\Psi.
\]
\end{proof}

\section{DLR intervals with a fixed particle number}\label{sec:blocks}

For a spatial sample let $a$ be the first positive point and let $b$ be
the $(K+2)$-nd positive point. The open interval $(a,b)$ contains exactly
$K$ points, written $x_1<\cdots<x_K$. Put
\[
L=K+1,\qquad \ell=b-a,\qquad c=L/\ell,\qquad
u_i=c(x_i-a).
\]
Let $\zeta=\gamma-\sum_{i=1}^K\delta_{x_i}$ denote the retained exterior.
Both endpoints, all negative points, and all points beyond $b$ are
retained. The conditioning sigma-field is
\[
\calG_K=\sigma(\zeta)
=\sigma\{\zeta(A):A\subset\R\text{ bounded Borel}\}.
\]
From $\zeta$, the point $a$ is recovered as its first positive point and
$b$ as its second positive point. On $\{b<R\}$ the original number of
points in $(0,R)$ is $\zeta((0,R))+K$. Thus the endpoints, the event
$\{b<R\}$, and this original particle number on that event are all
$\calG_K$-measurable.

\begin{samepage}
\begin{lemma}[DLR conditioning between particles of fixed ranks]\label{lem:random-dlr}
For every bounded measurable $f:\R^K\to\R$,
\begin{equation}\label{eq:random-dlr-law}
\E_P[f(x_1,\ldots,x_K)\mid\calG_K]
 =\int_{\Omega_{(a,b),K}}f(x)\,
        \kappa_{(a,b),K}^\beta(\zeta;\dd x)
 \quad P\text{-almost surely}.
\end{equation}
Let $S_{a,c}(x)=(c(x_1-a),\ldots,c(x_K-a))$ and let
$\widehat\zeta=(p\mapsto c(p-a))_\#\zeta$ be the scaled exterior.
The conditional distribution of the rescaled coordinates is
\begin{equation}\label{eq:rescaled-conditional-law}
\mu_\zeta=(S_{a,c})_\#\kappa_{(a,b),K}^\beta(\zeta;\cdot)
       =\kappa_{(0,L),K}^\beta(\widehat\zeta;\cdot),
\qquad
\E_P[f(u)\mid\calG_K]=\mu_\zeta(f).
\end{equation}
Its domain is $\Omega_K$, and $\widehat\zeta$ contains particles at
$0$ and $L$.
\end{lemma}
\end{samepage}

\begin{proof}
Fix a positive integer $R$ and work on $A_R=\{b<R\}$. Canonical DLR
in the deterministic interval $(0,R)$ first conditions on
$\gamma|_{(0,R)^c}$ and on $n=N_{(0,R)}$. For $n\ge K+2$, write its
ordered coordinates as $y_1<\cdots<y_n$. Further condition on
\[
y_1,\quad y_{K+2},y_{K+3},\ldots,y_n.
\]
The free coordinates are exactly $y_2,\ldots,y_{K+1}$, with domain
$y_1<y_2<\cdots<y_{K+1}<y_{K+2}$. In the deterministic-interval Gibbs
density, all terms involving only retained particles are constant
under this conditioning. The remaining terms are precisely the pair
interactions among these $K$ free particles and their interactions
with every retained particle, including $y_1=a$ and $y_{K+2}=b$.
Disintegration of the ordered Lebesgue density therefore gives the
normalized Gibbs kernel $\kappa_\zeta$ with exterior $\zeta$.

More explicitly, on $A_R$ the retained data in this two-stage
conditioning generate the same sigma-field as $\zeta$: the number $n$
is $\zeta((0,R))+K$, the retained coordinates are the points of
$\zeta$ in $(0,R)$, and the deterministic exterior is also read from
$\zeta$. For bounded measurable $h$ and $f$, disintegration on each
particle-number sector consequently gives
\[
\E_P\bigl[\1_{A_R}h(\zeta)f(x_1,\ldots,x_K)\bigr]
=\E_P\bigl[\1_{A_R}h(\zeta)\kappa_\zeta(f)\bigr].
\]
Every permitted interior resampling leaves $a$ the first positive point
and $b$ the $(K+2)$-nd positive point. The rank-selection indicator is
therefore identically one on this conditional domain; it contributes
no further factor to the density. Finally $A_R\uparrow1$ as integer
$R\to\infty$, so bounded convergence proves the displayed
conditional-expectation identity without its indicator. This identifies
the law given $\calG_K$.

The affine assertion follows from Lemma~\ref{lem:affine-covariance}
below. The Jacobian is constant and is absorbed by normalization.
\end{proof}

\begin{samepage}
\begin{lemma}[Affine covariance of relative exterior energies]
\label{lem:affine-covariance}
Let $I=(a,b)$ be a bounded interval with $K$ interior particles and a
locally finite exterior $\zeta$ of density one in both directions, for
which the symmetric relative exterior energies exist. For $c>0$, the
map $S_{a,c}:x\mapsto c(x-a)$ satisfies
\begin{equation}\label{eq:affine-kernel}
(S_{a,c})_\#\kappa_{I,K}^\beta(\zeta;\cdot)
 =\kappa_{(0,c(b-a)),K}^\beta
      ((p\mapsto c(p-a))_\#\zeta;\cdot).
\end{equation}
The same symmetric-cutoff convention is used on both sides. More
precisely, for the exterior summand $\Phi_p(x,y)$ defined in the proof,
\begin{equation}\label{eq:cutoff-center-error}
\lim_{S\to\infty}\sup_{x,y\in I^K}
\left|\sum_{\substack{p\in\zeta\\|p-a|\le S}}\Phi_p(x,y)
      -\sum_{\substack{p\in\zeta\\|p|\le S}}\Phi_p(x,y)\right|=0.
\end{equation}
\end{lemma}
\end{samepage}

\begin{proof}
For two interior configurations $x,y\in I^K$, the relative contribution
of an exterior point $p$ is, up to the fixed Hamiltonian sign,
\[
\Phi_p(x,y)=\sum_{i=1}^K
                    \log\frac{|x_i-p|}{|y_i-p|}.
\]
For $|p|\ge2(1+|a|+|b|)$ the mean value theorem gives
\[
\sup_{x,y\in I^K}|\Phi_p(x,y)|\le\frac{2K(b-a)}{|p|}.
\]
Compare cutoffs $|p|\le S$ and $|p-a|\le S$. Their symmetric
difference is contained in
\[
\mathcal S_S=\{p:S-|a|\le |p|\le S+|a|\}.
\]
Density one in each direction implies
$\zeta(\mathcal S_S)=o(S)$: the count on each fixed-width shell is
bounded by the difference of two cumulative counts of the form
$t+o(t)$. Hence
\[
\sup_{x,y\in I^K}
\left|\sum_{\substack{p\in\zeta\\|p-a|\le S}}\Phi_p(x,y)
     -\sum_{\substack{p\in\zeta\\|p|\le S}}\Phi_p(x,y)\right|
\le\frac{C K(b-a)}{S}\zeta(\mathcal S_S)=o(1).
\]
The uniformity here is for fixed $K,a,b$, exactly as needed for each
conditional law; no uniform rate in the random boundary is asserted.
Writing $q=c(p-a)$, $u_i=c(x_i-a)$ and $v_i=c(y_i-a)$ gives
\[
\log\frac{|u_i-q|}{|v_i-q|}
=\log\frac{|x_i-p|}{|y_i-p|}.
\]
A symmetric cutoff $|q|\le R$ corresponds to $|p-a|\le R/c$;
the preceding shell bound identifies its limit with the original
one. Relative pair energies are unchanged as well, since their
$\log c$ terms cancel between $x$ and $y$. Thus the transformed
relative Hamiltonian agrees with that of the scaled exterior. The
constant Jacobian $c^{-K}$ cancels in the normalized density. In
particular the endpoint particles are sent exactly to $0$ and $L$ in
the rescaling used above.
\end{proof}

Write $\gap_*=a-p_-$, where $p_-$ is the last point before the spatial
origin. This is the whole gap containing that origin. It is measurable
from $\zeta$, and $a\le\gap_*$.

\begin{samepage}
\begin{lemma}[Concentration of the endpoint distance]\label{lem:length}
The endpoint length satisfies
\begin{equation}\label{eq:length-probability}
\forall\varepsilon>0,\qquad
\lim_{K\to\infty}P\bigl(|\ell-L|>\varepsilon\sqrt K\bigr)=0.
\end{equation}
Equivalently, $\ell-L=o_P(\sqrt K)$. For any two candidate processes $P,Q$,
there is a deterministic sequence $r_K\downarrow0$ such that, for each
fixed $M$ and
\begin{equation}\label{eq:good}
E_{K,M}=\{\gap_*\le M,\ |\ell-L|\le r_K\sqrt K\},
\end{equation}
\[
\limsup_{K\to\infty}P(E_{K,M}^c)\le P(\gap_*>M),
\]
and the same assertion holds for $Q$.
The event $E_{K,M}$ is retained-exterior measurable.
\end{lemma}
\end{samepage}

\begin{proof}
Fix $\varepsilon>0$ and put
$t_\pm=K+2\pm\varepsilon\sqrt K$. For all sufficiently large $K$,
Chebyshev's inequality gives
\[
P(b>t_+)\le\frac{v_P(t_+)}{\varepsilon^2K}\longrightarrow0,
\qquad
P(b<t_-)\le\frac{v_P(t_-)}{\varepsilon^2K}\longrightarrow0,
\]
because $t_\pm/K\to1$ and \eqref{eq:hyper} holds. Now
\[
\frac{\ell-L}{\sqrt K}
 =\frac{b-(K+2)}{\sqrt K}+\frac{1-a}{\sqrt K}
 \longrightarrow0\quad\text{in }P\text{-probability}.
\]
The variable $a$ is finite almost surely and does not depend on $K$.
A diagonal choice gives a deterministic $r_K\downarrow0$ such that
\[
P(|\ell-L|>r_K\sqrt K)+Q(|\ell-L|>r_K\sqrt K)\longrightarrow0.
\]
The assertions about $E_{K,M}$ follow.
\end{proof}

Define the conditional quadratic cost
\begin{equation}\label{eq:J}
J_K(\zeta)=\E_{\mu_\zeta}\sum_{i=1}^K(u_i-i)^2.
\end{equation}

\begin{samepage}
\begin{lemma}[Subquadratic expected displacement]\label{lem:quadratic}
For every fixed $M$,
\begin{equation}\label{eq:J-small}
K^{-2}\E_P[\1_{E_{K,M}}J_K]\longrightarrow0.
\end{equation}
\end{lemma}
\end{samepage}

\begin{proof}
For deterministic ordered $z_1,\ldots,z_K\in(0,\ell)$, let $A(t)$ be
their counting function and $\rho=K/\ell$. Direct expansion yields the
exact identity
\begin{equation}\label{eq:quantile}
\int_0^\ell(A(t)-\rho t)^2\dd t
=\rho\sum_{i=1}^K
\left(z_i-\frac{i-\frac12}{\rho}\right)^2+\frac{\ell}{12}.
\end{equation}
For example,
$\int_0^\ell A(t)^2\dd t=\sum_i(2i-1)(\ell-z_i)$ and
$\int_0^\ell tA(t)\dd t=\frac12\sum_i(\ell^2-z_i^2)$, which verify
\eqref{eq:quantile}.

Take $z_i=x_i-a$ and use $D_0$ from \eqref{eq:oriented-discrepancy}. Except at endpoints of
Lebesgue measure zero,
\[
A(t)-\rho t=D_0(a+t)+(a-1)+(1-\rho)t.
\]
On $E_{K,M}$, $a\le M$, $\ell\asymp K$, and $b\le3K$ for all
sufficiently large $K$ (with $M$ fixed). To compare the two deterministic particle locations, set
$m_i=(i-\frac12)\ell/K$. Since $L=K+1$,
\[
\left|m_i-\frac{i\ell}{L}\right|
 =\frac\ell K\left|\frac iL-\frac12\right|
 \le\frac\ell{2K}\qquad(1\le i\le K).
\]
Consequently, \eqref{eq:quantile} yields
\[
\sum_{i=1}^K(u_i-i)^2
\le\frac{2c^2}{\rho}\int_0^\ell(A(t)-\rho t)^2\dd t
   +\frac{L^2}{2K}.
\]
Also,
\[
\int_0^\ell(A(t)-\rho t)^2\dd t
\le3\int_0^{3K}D_0(s)^2\dd s
  +3\ell(a-1)^2+\ell(\ell-K)^2.
\]
Here $c^2/\rho=L^2/(K\ell)$ is bounded on $E_{K,M}$, and
$(\ell-K)^2\le2(\ell-L)^2+2$. These inequalities give, on this event,
\begin{equation}\label{eq:quadratic-bound}
\sum_i(u_i-i)^2
\le C\int_0^{3K}D_0(s)^2\dd s+C_MK+CK(\ell-L)^2.
\end{equation}
By \eqref{eq:hyper},
$\int_0^{3K}v_P(s)\dd s=o(K^2)$, and on $E_{K,M}$ the last term
in \eqref{eq:quadratic-bound} is at most $Cr_K^2K^2$. Therefore
\[
\frac1{K^2}\E_P\!\left[\1_{E_{K,M}}
     \sum_{i=1}^K(u_i-i)^2\right]
\le\frac C{K^2}\int_0^{3K}v_P(s)\dd s+\frac{C_M}{K}+Cr_K^2
\longrightarrow0.
\]
Because the good event is retained-exterior measurable, conditional
expectation turns the left side into the expression in
\eqref{eq:J-small}.
\end{proof}

\section{Tightness of the conditional exterior-force contribution}\label{sec:forces}

The force estimates will only use tightness. We will not require an
integrable product of an exterior energy and an interior count envelope.

\Needspace{13\baselineskip}
\begin{samepage}
\begin{lemma}[A square-integrable bound for the half-line force]\label{lem:halfline}
Let $q$ be a particle and let the exterior lie on one side of $q$.
After deleting $q$, define its outward counting function by one of
\[
N_{\rm ext}(t)=\gamma([q-t,q))
\quad\text{or}\quad
N_{\rm ext}(t)=\gamma((q,q+t]),\qquad t\ge0.
\]
Assume the configuration belongs to the regularity set of
Lemma~\ref{lem:as-regularity}. Define
\begin{equation}\label{eq:halfline-definitions}
D(t)=N_{\rm ext}(t)-(t-1)_+,\qquad
H(x)=\int_0^\infty\frac{|D(t)|}{(x+t)^2}\dd t,\quad x\ge0.
\end{equation}
The function $(t-1)_+$ is the cumulative mass of the continuum ray
$[1,\infty)$ with unit density.
The associated bound satisfies
\begin{equation}\label{eq:B}
\calB=H(0)^2+\int_0^\infty H(x)^2\dd x<\infty.
\end{equation}
For $c\in[1/2,2]$, define the rescaled quantities by
\begin{equation}\label{eq:halfline-scaled-definitions}
\begin{aligned}
D^{(c)}(t)&=N_{\rm ext}(t/c)-c^{-1}(t-1)_+,\\
H^{(c)}(x)&=\int_0^\infty\frac{|D^{(c)}(t)|}{(x+t)^2}\dd t,\\
\calB^{(c)}&=H^{(c)}(0)^2+\int_0^\infty H^{(c)}(x)^2\dd x.
\end{aligned}
\end{equation}
There is an absolute constant $C<\infty$ such that
\begin{equation}\label{eq:halfline-scale-bound}
\sup_{1/2\le c\le2}\calB^{(c)}\le C(\calB+1).
\end{equation}
\end{lemma}
\end{samepage}

\begin{proof}
Simplicity and local finiteness give $d\in(0,1]$ such that no
residual exterior point has outward distance less than $d$. Thus
$D(t)=0$ on $[0,d)$. Lemma~\ref{lem:as-regularity} controls the tail,
and Cauchy--Schwarz gives
\[
\int_0^\infty\frac{D(t)^2}{t^2}\dd t<\infty,
\qquad
H(0)^2\le\frac1d\int_0^\infty\frac{D(t)^2}{t^2}\dd t<\infty.
\]
For completeness, the Schur-test kernel $K(x,t)=t/(x+t)^2$ satisfies
\[
\int_0^\infty K(x,t)t^{-1/2}\dd t=\frac\pi2x^{-1/2},\qquad
\int_0^\infty K(x,t)x^{-1/2}\dd x=\frac\pi2t^{-1/2}.
\]
Since $H(x)=\int_0^\infty K(x,t)|D(t)|/t\dd t$, the Schur test gives
\begin{equation}\label{eq:hardy}
\int_0^\infty H(x)^2\dd x
\le\frac{\pi^2}{4}\int_0^\infty D(t)^2t^{-2}\dd t.
\end{equation}

For the scaled configuration the centered discrepancy is
\[
D(t/c)+(t/c-1)_+-c^{-1}(t-1)_+.
\]
Write the last two terms as
$R_c(t)=(t/c-1)_+-c^{-1}(t-1)_+$. For $c\in[1/2,2]$,
\[
|R_c(t)|\le C\1_{\{t\ge1/2\}},\qquad
H^{(c)}(x)\le c^{-1}H(x/c)+\frac C{x+1/2}.
\]
Squaring, integrating in $x$, and evaluating at $x=0$ proves
\eqref{eq:halfline-scale-bound}, with the same constant for every
$c\in[1/2,2]$.
\end{proof}

Let $B_\zeta(u)$ be the difference of the actual and lattice exterior
forces in the rescaled block. More precisely, with endpoint terms
cancelled,
\[
B_\zeta(u)=\PV\sum_{p\in\widehat\zeta\setminus\{0,L\}}\frac1{u-p}
-\PV\sum_{p\in\xi_K^{\rm lat}\setminus\{0,L\}}\frac1{u-p},
\qquad 0\le u\le L.
\]
For almost every retained exterior under a candidate law, each paired
sum exists and their difference is bounded on $[0,L]$. Define
\begin{equation}\label{eq:Dcost}
\mathcal D_K(\zeta)=
\E_{\mu_\zeta}\sum_{i=1}^K B_\zeta(u_i)^2.
\end{equation}

\begin{samepage}
\begin{proposition}[Tightness of the conditional squared force difference]\label{prop:force}
For fixed $M$, the nonnegative random variables
$\1_{E_{K,M}}\mathcal D_K$ are tight as $K\to\infty$. Explicitly,
\begin{equation}\label{eq:force-tight-tail}
\lim_{A\to\infty}\limsup_{K\to\infty}
P\bigl(E_{K,M}\cap\{\mathcal D_K>A\}\bigr)=0.
\end{equation}
The statement holds separately for each candidate process, with no
assertion of a uniform first moment for $\mathcal D_K$.
\end{proposition}
\end{samepage}

\begin{proof}
On $E_{K,M}$, $c=L/\ell\in[1/2,2]$ for all sufficiently large $K$.
The scaled exterior $\widehat\zeta$ has density
$\lambda=c^{-1}=\ell/L$. Define its left and right outward counts,
centered discrepancies, and nonnegative bounds by
\begin{equation}\label{eq:boundary-halfline-quantities}
\begin{aligned}
N_{\rm L}(t)&=\widehat\zeta([-t,0)),&
N_{\rm R}(t)&=\widehat\zeta((L,L+t]),\\
D_\star(t)&=N_\star(t)-\lambda(t-1)_+,&
H_\star(x)&=\int_0^\infty\frac{|D_\star(t)|}{(x+t)^2}\dd t,\\
\calB_\star&=H_\star(0)^2+\int_0^\infty H_\star(x)^2\dd x,
& &\star\in\{{\rm L},{\rm R}\}.
\end{aligned}
\end{equation}
Both endpoint particles have already been removed from these counts.
For the residual unit lattice, set
\[
D_{\rm lat}(t)=\lfloor t\rfloor-(t-1)_+,\qquad
0\le D_{\rm lat}(t)\le\1_{\{t\ge1\}}.
\]
Stieltjes integration by parts, first with finite cutoffs, gives the
exact signed decomposition
\begin{equation}\label{eq:force-decomposition}
\begin{split}
B_\zeta(u)
={}&\int_0^\infty\frac{D_{\rm L}(t)-D_{\rm lat}(t)}{(u+t)^2}\dd t
   -\int_0^\infty\frac{D_{\rm R}(t)-D_{\rm lat}(t)}{(L-u+t)^2}\dd t\\
&+(\lambda-1)\log\frac{L-u+1}{u+1}.
\end{split}
\end{equation}
Indeed, the boundary terms at infinity vanish because
$D_\star(t)=o(t)$, and the continuum contribution is
\[
\int_1^\infty\left(\frac1{u+t}-\frac1{L-u+t}\right)\dd t
 =\log\frac{L-u+1}{u+1}.
\]
The change from a common outward cutoff to the symmetric spatial
cutoff has vanishing error, by the same fixed-width shell estimate as
in Lemma~\ref{lem:affine-covariance}. Finally,
\[
\int_0^\infty\frac{D_{\rm lat}(t)}{(x+t)^2}\dd t
\le\frac1{x+1},\qquad x\ge0.
\]
Consequently,
\begin{equation}\label{eq:force-pointwise}
\begin{split}
|B_\zeta(u)|\le{}&
H_{\rm L}(u)+H_{\rm R}(L-u)+\frac C{u+1}+\frac C{L-u+1}\\
&+|\lambda-1|\abs{\log\frac{u+1}{L-u+1}}.
\end{split}
\end{equation}
The endpoint charges at $0,L$ cancel exactly before this estimate.
The quantities $\calB_{\rm L}$ and $\calB_{\rm R}$ in
\eqref{eq:boundary-halfline-quantities} depend only on the retained exterior.

For an interior configuration $u\in\Omega_K$, define
\begin{equation}\label{eq:interior-count-definitions}
C_{\rm int}^{\rm L}(u)=\sup_{t\ge0}
 \frac{\#\{i:u_i\le t\}}{1+t},\qquad
C_{\rm int}^{\rm R}(u)=\sup_{t\ge0}
 \frac{\#\{i:L-u_i\le t\}}{1+t}.
\end{equation}
For any nonnegative nonincreasing function $h:[0,L]\to[0,\infty)$,
Stieltjes integration by parts gives
\begin{equation}\label{eq:sum-envelope}
\sum_i h(u_i)^2
\le C_{\rm int}^{\rm L}
\left(h(0)^2+\int_0^L h(s)^2\dd s\right).
\end{equation}
For instance, set $N_{\rm int}(t)=\#\{i:u_i\le t\}$ and use
$N_{\rm int}(t)\le C_{\rm int}^{\rm L}(1+t)$ in
$\int h^2\dd N_{\rm int}$; the terminal boundary term gives exactly
the expression on the right. The right-endpoint version is
\[
\sum_{i=1}^K h(L-u_i)^2
\le C_{\rm int}^{\rm R}
\left(h(0)^2+\int_0^Lh(s)^2\dd s\right).
\]

For $h_L(s)=\log((L+1)/(s+1))$, one has
\[
\left|\log\frac{u+1}{L-u+1}\right|
 \le h_L(u)+h_L(L-u),\qquad
h_L(0)^2+\int_0^Lh_L(s)^2\dd s\le CL.
\]
Squaring \eqref{eq:force-pointwise}, using
\eqref{eq:sum-envelope}, and taking conditional expectations gives
\begin{equation}\label{eq:force-total}
\mathcal D_K\le
C\{X_{\rm L}(\calB_{\rm L}+1)+X_{\rm R}(\calB_{\rm R}+1)\}
+C(\lambda-1)^2L(X_{\rm L}+X_{\rm R}),
\end{equation}
where the exterior-measurable conditional expectations are
\begin{equation}\label{eq:conditional-count-definitions}
X_\star(\zeta)=\mu_\zeta(C_{\rm int}^\star),
\qquad \star\in\{{\rm L},{\rm R}\}.
\end{equation}

We verify \eqref{eq:force-tight-tail} under the distribution of the retained
exterior induced by $P$, in four steps. All bounds below are for fixed $M$ and sufficiently large $K$.

\emph{Exterior envelopes.}
Let $\calB^{\rm out,-}$ and $\calB^{\rm out,+}$ be the unscaled
leftward and rightward quantities of Lemma~\ref{lem:halfline}, viewed
as functions of a Palm-rooted configuration. They are finite
$P^0$-almost surely, but no expectation bound for them is needed.
The actual left endpoint is $a=a_0$, and the right endpoint is
$b=a_{K+1}$. The scale estimate in that lemma implies on $E_{K,M}$
\[
\calB_{\rm L}\le C\{\calB^{\rm out,-}(\theta_a\gamma)+1\},\qquad
\calB_{\rm R}\le C\{\calB^{\rm out,+}(\theta_b\gamma)+1\}.
\]
Both right sides depend only on the appropriate outward half-line.
Using Lemma~\ref{lem:endpoint-domination}, uniformly in $K$,
\begin{align*}
P(E_{K,M}\cap\{\calB_{\rm L}>A\})
&\le M P^0(\calB^{\rm out,-}>A/C-1),\\
P(E_{K,M}\cap\{\calB_{\rm R}>A\})
&\le M P^0(\calB^{\rm out,+}>A/C-1).
\end{align*}
The right sides tend to zero as $A\to\infty$. The domination constant
at the right endpoint is unchanged although its point-shift index
$K+1$ grows with $K$.

\emph{Conditional count envelopes.}
For an actual sample on the good event, $c\in[1/2,2]$ gives
\[
C_{\rm int}^{\rm L}\le2C^+(\theta_a\gamma),\qquad
C_{\rm int}^{\rm R}\le2C^-(\theta_b\gamma).
\]
For example $\#\{i:u_i\le t\}\le
N_{(a,a+t/c]}(\gamma)\le C^+(\theta_a\gamma)(1+t/c)$,
and $(1+t/c)/(1+t)\le2$. By Lemma~\ref{lem:random-dlr},
$X_\star=\E_P[C_{\rm int}^\star\mid\calG_K]$. Since $E_{K,M}$ belongs
to $\calG_K$, Lemmas~\ref{lem:endpoint-domination} and
\ref{lem:palm-envelope} imply
\begin{equation}\label{eq:conditional-count-moment}
\begin{aligned}
\E_P[\1_{E_{K,M}}X_{\rm L}]
 &=\E_P[\1_{E_{K,M}}C_{\rm int}^{\rm L}]
 \le2M\E_{P^0}C^+,\\
\E_P[\1_{E_{K,M}}X_{\rm R}]
 &=\E_P[\1_{E_{K,M}}C_{\rm int}^{\rm R}]
 \le2M\E_{P^0}C^-.
\end{aligned}
\end{equation}
The bound is uniform in $K$. Markov's inequality therefore makes
$\1_{E_{K,M}}X_\star$ uniformly tight.

\emph{Products.}
For $\star\in\{{\rm L},{\rm R}\}$ and any $A>0$, the union bound gives
\begin{align*}
P\bigl(E_{K,M}\cap\{X_\star(\calB_\star+1)>A\}\bigr)
&\le P(E_{K,M}\cap\{X_\star>\sqrt A\})\\
&\quad+P(E_{K,M}\cap\{\calB_\star+1>\sqrt A\}).
\end{align*}
The limit superior as $K\to\infty$ of each term tends to zero as
$A\to\infty$. This
proves tightness of the products using their marginal tails only;
neither independence nor an integrable product has been used.

\emph{Density mismatch.}
On $E_{K,M}$,
\[
(\lambda-1)^2L=\frac{(\ell-L)^2}{L}\le Cr_K^2.
\]
Consequently \eqref{eq:conditional-count-moment} gives the bound
\[
\E_P\bigl[\1_{E_{K,M}}(\lambda-1)^2L(X_{\rm L}+X_{\rm R})\bigr]
\le C_Mr_K^2\longrightarrow0.
\]
This particular term vanishes in mean. Combining its vanishing tail
with the two tight product terms in \eqref{eq:force-total} proves
\eqref{eq:force-tight-tail}.
\end{proof}

\section{Proof of the discrepancy comparison theorem}
\label{sec:proof-main}

Let $P,Q$ satisfy Theorem~\ref{thm:comparison}. Fix $r,G$ and first fix
$M<\infty$. Choose $r_K$ in Lemma~\ref{lem:length} jointly for $P$ and
$Q$, and use this same sequence in both good events $E_{K,M}$.

\subsection{Cancellation in the logarithmic derivative}

For almost every retained exterior, differentiating the two Gibbs
densities gives, for $1\le i\le K$,
\begin{equation}\label{eq:score-identity}
\partial_i\log\frac{\dd\mu_\zeta}{\dd\omega_{K,T}}(u)
=\beta B_\zeta(u_i)+\frac{u_i-i}{T}.
\end{equation}
Indeed, the interior pair Hamiltonian is identical in the two laws,
and the endpoint charges at $0,L$ also agree. Both cancel in the
density ratio. The remaining exterior-force difference is exactly
$B_\zeta$, with the convention of Section~\ref{sec:forces}.
The inequality $(a+b)^2\le2a^2+2b^2$ therefore gives
\begin{equation}\label{eq:score}
I(\mu_\zeta\mid\omega_{K,T})
\le 2\beta^2\mathcal D_K(\zeta)+\frac{2J_K(\zeta)}{T^2}.
\end{equation}
For each fixed exterior the residual force is bounded on $[0,L]$:
the adjacent exterior particles have positive distance from the
endpoints, and the discrepancy envelope controls the tails.
Thus the logarithm of the density ratio extends to a bounded Lipschitz
function on the simplex closure, and its entropy and Fisher information
are finite. This cancellation avoids inverse-gap moment assumptions,
including when $0<\beta<1$.

\subsection{A common confinement scale}

Define
\[
q_{K,M}=\frac1{K^2}\left\{
\E_P[\1_{E_{K,M}}J_K]+\E_Q[\1_{E_{K,M}}J_K]\right\}\longrightarrow0,
\]
where the convergence is Lemma~\ref{lem:quadratic} for the fixed $M$.
Set
\begin{equation}\label{eq:choice}
\varepsilon_{K,M}=\max\{\sqrt{q_{K,M}},K^{-1/2}\},
\qquad T_{K,M}=\varepsilon_{K,M}K.
\end{equation}
These deterministic sequences may depend on $P,Q,M$. The same
reference $\omega_{K,T_{K,M}}$ is used in both comparisons. In
particular,
\begin{equation}\label{eq:scale-consequences}
\varepsilon_{K,M}\to0,\qquad
T_{K,M}\ge\sqrt K\to\infty,\qquad
\frac{T_{K,M}}K\to0,\qquad
\frac{q_{K,M}}{\varepsilon_{K,M}}
\le\sqrt{q_{K,M}}\to0.
\end{equation}
Thus the confinement becomes weak, while its scale is still small
compared with the block length.

Proposition~\ref{prop:gap-entropy} and \eqref{eq:score} imply
\begin{equation}\label{eq:comparison-final}
|\mu_\zeta(F_K)-\omega_{K,T_{K,M}}(F_K)|^2
\le C_G\left\{\varepsilon_{K,M}\mathcal D_K
                 +\frac{J_K}{\varepsilon_{K,M}K^2}\right\}.
\end{equation}
The two terms are controlled separately. For either candidate law
$R\in\{P,Q\}$, Proposition~\ref{prop:force} and
$\varepsilon_{K,M}\to0$ imply
\begin{equation}\label{eq:force-vanishing-probability}
\forall\delta>0,\qquad
R\!\left(E_{K,M}\cap
\{\varepsilon_{K,M}\mathcal D_K>\delta\}\right)\longrightarrow0.
\end{equation}
To see the use of tightness explicitly, for every fixed $A>0$ and all
large $K$ with $\varepsilon_{K,M}A\le\delta$, this probability is
at most $R(E_{K,M}\cap\{\mathcal D_K>A\})$. Take first the limit
superior in $K$ and then let $A\to\infty$. For the quadratic term, under $P$,
\[
\E_P\left[\frac{\1_{E_{K,M}}J_K}
{\varepsilon_{K,M}K^2}\right]
\le\frac{q_{K,M}}{\varepsilon_{K,M}}\longrightarrow0,
\]
and the same bound holds under $Q$.
Consequently, for $R\in\{P,Q\}$,
\begin{equation}\label{eq:conditional-comparison-L1}
\E_R\!\left[\1_{E_{K,M}}
 \left|\mu_\zeta(F_K)-\omega_{K,T_{K,M}}(F_K)\right|\right]
 \longrightarrow0.
\end{equation}
Indeed, the integrand converges to zero in $R$-probability by
\eqref{eq:comparison-final} and is bounded by $2\norm G_\infty$.
This deduction does not use a first-moment bound on $\mathcal D_K$.

\subsection{Order of limits and identification}

Let $\mathfrak F_K(\gamma)$ be the rescaled block statistic
$F_K(u(\gamma))$. Conditional DLR and exterior measurability of
$E_{K,M}$ imply
\[
\limsup_{K\to\infty}
\abs{\E_P\mathfrak F_K-\omega_{K,T_{K,M}}(F_K)}
\le2\norm G_\infty P(\gap_*>M).
\]
The analogous estimate holds for $Q$. For this fixed $M$, cancelling
the common reference gives
\[
\limsup_{K\to\infty}
\abs{\E_P\mathfrak F_K-\E_Q\mathfrak F_K}
\le2\norm G_\infty\{P(\gap_*>M)+Q(\gap_*>M)\}.
\]
Only now let $M\to\infty$. The gap $\gap_*$ is finite almost surely,
so
\begin{equation}\label{eq:cancel}
\lim_{K\to\infty}
\abs{\E_P\mathfrak F_K-\E_Q\mathfrak F_K}=0.
\end{equation}
No estimate uniform in $M$ is needed: the common reference expectation
has been eliminated for each fixed $M$ before taking $M\to\infty$.

The affine rescaling does not alter the limiting statistic. Extending
$G$ by zero gives a bounded globally Lipschitz function, and each gap
occurs in at most $r$ windows. The unscaled statistic is $\mathfrak F_K^{\rm un}=A_{K,r}$ from
\eqref{eq:spatial-gap-average}, because $x_i=a_i$. Thus
\[
|\mathfrak F_K-\mathfrak F_K^{\rm un}|
\le\frac{C_G}{K-r}|c-1|\ell
=\frac{C_G}{K-r}|\ell-L|\longrightarrow0
\]
almost surely by density one, and also in mean by boundedness.
Lemma~\ref{lem:identify} and \eqref{eq:cancel} yield
\[
\E_{P^0}G(g_0,\ldots,g_{r-1})
=\E_{Q^0}G(g_0,\ldots,g_{r-1})
\]
for every $r$ and $G\in C_c^2((0,\infty)^r)$.
These test functions determine probability measures on
$(0,\infty)^r$, so every finite consecutive-gap distribution agrees.
Point-shift invariance moves any finite integer-indexed gap window
into such a consecutive window. It follows that
\begin{equation}\label{eq:gap-law-equality}
\Law_{P^0}\bigl((g_j)_{j\in\Z}\bigr)
 =\Law_{Q^0}\bigl((g_j)_{j\in\Z}\bigr).
\end{equation}

The Palm configuration is reconstructed from its gaps and the point
at zero by
\[
p_n=\sum_{j=0}^{n-1}g_j,\qquad
p_{-n}=-\sum_{j=-n}^{-1}g_j,\qquad n\ge1.
\]
Thus $P^0=Q^0$. Finally Palm inversion at intensity one gives, for
bounded measurable $f$,
\[
\E_P f(\gamma)=
\E_{P^0}\left[\int_0^{g_0}f(\theta_t\gamma)\dd t\right],
\]
and the same formula for $Q$ proves $P=Q$. This completes
Theorem~\ref{thm:comparison}, without an ergodicity assumption.


\section{From logarithmic electric energy to discrepancy control}
\label{sec:energy}

Theorem~\ref{thm:comparison} has been proved without an electric-energy
assumption. This section verifies its three variance hypotheses for
the ergodic finite-energy components used in
Theorem~\ref{thm:finite-energy}. We separate the imported linear
discrepancy estimate from the integrated-variance conclusion proved
below.

\subsection{Energy convention and results used from the literature}

For a configuration $\gamma\in\Conf(\R)$, distinguish its embedding
in the plane from the one-dimensional counting measure:
\[
\overline\gamma=\sum_{q\in\gamma}\delta_{(q,0)},\qquad
\dd x\,\delta_0=\Leb\otimes\delta_0.
\]
The second measure is the unit-density background on the horizontal
axis. Integrals over subsets of $\R^2$ use $\dd x\dd y$ unless a
different measure is displayed. The notation $\langle T,\phi\rangle$
denotes the action of a distribution $T$ on a test function $\phi$.
Fix $p\in(1,2)$ and define
\begin{equation}\label{eq:compatible-fields}
\operatorname{Comp}(\gamma)=
\left\{E\in L^p_{\rm loc}(\R^2;\R^2):
 -\diver E=2\pi(\overline\gamma-\dd x\,\delta_0)
 \text{ in distributions}\right\}.
\end{equation}
Thus the divergence constraint means
\[
\int_{\R^2}E\cdot\nabla\phi
 =2\pi\left\{\sum_{q\in\gamma}\phi(q,0)
                  -\int_\R\phi(x,0)\dd x\right\}
\qquad(\phi\in C_c^\infty(\R^2)).
\]
No curl-free condition is imposed in this definition.

For $0<\eta<1$, let $\rho_\eta$ be the uniform probability measure
on the circle of radius $\eta$, characterized by
\[
\int\phi(z)\rho_\eta(\dd z)
 =\frac1{2\pi}\int_0^{2\pi}
   \phi(\eta\cos\theta,\eta\sin\theta)\dd\theta.
\]
Define, almost everywhere in $z\in\R^2$,
\begin{equation}\label{eq:inner-truncation}
f_\eta(z)=\1_{\{|z|\le\eta\}}\log(|z|/\eta),\qquad
E_\eta=E+\sum_{q\in\gamma}\nabla f_\eta(\,\cdot-(q,0)).
\end{equation}
The value of $f_\eta$ at zero is irrelevant to this locally integrable
function. Since $\Delta f_\eta=2\pi(\delta_{(0,0)}-\rho_\eta)$,
the truncated field $F=E_\eta$ satisfies
\begin{equation}\label{eq:regularized-div}
-\diver F=2\pi(\overline\gamma*\rho_\eta-dx\,\delta_0).
\end{equation}
Define
\begin{equation}\label{eq:renormalized-energy}
\begin{aligned}
\mathcal W_\eta(E)
&=\limsup_{R\to\infty}\frac1{2R}
\int_{[-R,R]\times\R}|E_\eta(x,y)|^2\dd x\dd y
+2\pi\log\eta,\\
\mathcal W(E)&=\lim_{\eta\downarrow0}\mathcal W_\eta(E),\\
W(\gamma)&=\inf_{E\in\operatorname{Comp}(\gamma)}\mathcal W(E).
\end{aligned}
\end{equation}
The spatial average is per unit horizontal length. The order of operations
is $R\to\infty$, then $\eta\downarrow0$, and finally the infimum
over compatible fields. The limits have the extended-real interpretation
of the regularized energy, and $\inf\varnothing=+\infty$. The truncation-limit existence is the one recorded in
\cite[Section~2.7.1]{LS}.

The choice of $p\in(1,2)$ does not change the finite-energy class.
Indeed, if a locally integrable compatible field has finite
$\mathcal W(E)$, some fixed $\eta$ has finite spatial mean square
energy. Then $E_\eta\in L^2_{\rm loc}$, and on each compact set
$E-E_\eta$ is a finite sum of truncated Coulomb singularities. Each
such gradient belongs to $L^p_{\rm loc}$ for every $p<2$. Hence every
field relevant to a finite infimum in \eqref{eq:renormalized-energy}
belongs to the displayed admissible class.

\begin{samepage}
\begin{lemma}[Energy normalization and the linear variance bound]\label{lem:energy-input}
Write $W_{\rm LS}(\gamma,1)$ and $\widetilde W_{\rm LS}(P,1)$ for the
configuration and process energies of \cite[Sections~2.7.3--2.7.4]{LS}
in the Log1 case. Then
\begin{equation}\label{eq:energy-bridge}
W(\gamma)=2\pi W_{\rm LS}(\gamma,1),\qquad
\E_PW=2\pi\widetilde W_{\rm LS}(P,1).
\end{equation}
Writing $\widetilde{\mathcal W}^{\rm div}$ and
$\mathcal W^{\rm div}$ for the configuration and process energies
of \cite[Definitions~2.5--2.6]{Assiotis}, one also has
\begin{equation}\label{eq:assiotis-energy}
W(\gamma)=2\pi\widetilde{\mathcal W}^{\rm div}(\gamma),\qquad
\E_PW=2\pi\mathcal W^{\rm div}(P).
\end{equation}
Thus the finite-energy hypothesis in Theorem~\ref{thm:finite-energy}
is exactly the energy hypothesis of \cite[Theorem~1.1]{Assiotis}.
In particular, there is an absolute constant $C_0<\infty$ such that
$W(\gamma)\ge-C_0$. If $P$ is stationary and $\E_PW<\infty$, then
\begin{equation}\label{eq:energy-linear-conclusion}
\E_PN_A=\Leb(A),\qquad
v_P(t)\le C_Pt\quad(t>0),
\end{equation}
for every bounded Borel set $A$, with $C_P<\infty$ depending on $P$.
\end{lemma}
\end{samepage}

\begin{proof}
Specialize \cite[Sections~2.6.4--2.7.4]{LS} to $d=1$, extension
dimension $k=1$, weight exponent zero, background $m=1$, and
$c_{d,s}=2\pi$. The admissible field class and divergence equation
then agree with ours. Their positive cutoff is
\[
f_\eta^{\rm LS}(z)=(-\log|z|+\log\eta)_+=-f_\eta(z),
\]
and their truncation $E-\sum\nabla f_\eta^{\rm LS}$ is exactly
$E_\eta$. Their cube $\square_S$ has side length $S$; taking $S=2R$
identifies their spatial average with ours. Finally their subtraction
$-c_{d,s}g(\eta)$ is $+2\pi\log\eta$, and their configuration energy
includes a factor $1/c_{d,s}$. This proves \eqref{eq:energy-bridge}.
The lower bound is the standard one recalled in the proof of
\cite[Lemma~3.9]{LS}.

For \eqref{eq:assiotis-energy}, take the exponent $q_{\rm el}=p$ in
\cite[Definition~2.5]{Assiotis}. Its field class has exactly the same
divergence constraint, with no curl-free requirement. Its positive
inner cutoff is $f_\eta^{\rm LS}$, so the regularized fields agree.
Definition~2.6 there uses the same strip, order of limits, and infimum,
with regularized functional
\[
\frac1{2R\,2\pi}\int_{[-R,R]\times\R}|E_\eta|^2
+\log\eta.
\]
This is our entire regularized functional divided by $2\pi$.
Positive scaling commutes with the limits and infimum, giving the
configuration identity. The process energy in that definition is the
expectation of this configuration infimum, so the process identity
follows directly; no optimization over stationary joint distributions of configurations and
fields, and no measurable field selection, is involved.

By \cite[Lemma~3.2]{LS} and \eqref{eq:energy-bridge}, for $t\ge1$,
\[
\E_P(N_{(0,t]}-t)^2\le C(C+\E_PW)t.
\]
Stationarity forces intensity one, since any different intensity would
give a quadratic lower bound by Jensen's inequality.
For completeness, the small-interval bound follows from stationarity
alone once $\E_PN_{(0,1]}^2<\infty$. If $0<t<1$ and
$m=\lfloor1/t\rfloor$, then
\[
m\E_PN_{(0,t]}^2
=\E_P\sum_{j=0}^{m-1}N_{(jt,(j+1)t]}^2
\le\E_PN_{(0,1]}^2.
\]
Since $m\ge1/(2t)$, this gives
$v_P(t)\le2t\E_PN_{(0,1]}^2$ and completes \eqref{eq:linear}.
\end{proof}

\begin{remark}[Truncation versus normalization]\label{rem:cutoff}
The cutoff displayed in \cite[Definition~2.20]{DHLM}, in the
arXiv version specified in the bibliography, has $|z|\ge\eta$.
We use the compatible underlying formulation of \cite{LS}, matched
explicitly in Lemma~\ref{lem:energy-input}.
The singular energy of a unit charge is
\[
\int_{\eta<|z|<r_0}|\nabla\log|z||^2\dd z
=2\pi\log(r_0/\eta).
\]
Thus $1/(2R)$ over $[-R,R]\times\R$ pairs with $+2\pi\log\eta$.
Using $1/R$ over that window would require $+4\pi\log\eta$.
Multiplying the entire regularized functional, including its
subtraction, by a positive constant preserves finiteness. Changing
the cutoff or only one of those two coefficients has no such
justification. The identity \eqref{eq:energy-bridge} documents the
applicability of the imported energy results.
\end{remark}

\subsection{A stationary configuration--field law and the spectral estimate}


For vector fields, the horizontal translation is
$(\theta_aF)(x,y)=F(x+a,y)$. A joint law of $(\gamma,F)$ is called
stationary when it is invariant under
$(\gamma,F)\mapsto(\theta_a\gamma,\theta_aF)$ for every $a\in\R$.

\begin{samepage}
\begin{lemma}[Stationary joint law of a configuration and a field]\label{lem:lift}
Let $P$ be ergodic and stationary, and suppose
$P(W<\infty)=1$, where $W$ is defined in
\eqref{eq:renormalized-energy}. There exist $\eta\in(0,1)$ and a
probability measure $\mathsf Q$ on configuration--field pairs such that
\begin{equation}\label{eq:joint-law-properties}
\begin{aligned}
\Law_{\mathsf Q}(\gamma)&=P,\\
(\theta_a,\theta_a)_\#\mathsf Q&=\mathsf Q\quad(a\in\R),\\
-\diver F&=2\pi(\overline\gamma*\rho_\eta-\dd x\,\delta_0)
\quad\mathsf Q\text{-almost surely},
\end{aligned}
\end{equation}
and
\begin{equation}\label{eq:field-energy}
\calE=\E_{\mathsf Q}\int_{[0,1]\times\R}|F(x,y)|^2\dd x\dd y<\infty.
\end{equation}
\end{lemma}
\end{samepage}

\begin{proof}
Choose a configuration $\gamma$ with $W(\gamma)<\infty$ such that
\begin{equation}\label{eq:generic-configuration}
\lim_{R\to\infty}\frac1{2R}\int_{-R}^R
 f_j(\theta_x\gamma)\dd x=P(f_j),\qquad j\ge1,
\end{equation}
where $(f_j)_{j\ge1}$ is a countable convergence-determining family of
bounded continuous functions on the space of locally finite measures
with the vague topology. Ergodicity ensures that
such configurations have full probability. Choose a compatible field
whose renormalized functional is finite; attainment of the infimum
defining $W$ is unnecessary. At some sufficiently small fixed $\eta$,
the truncated field $F=E_\eta$ consequently satisfies
\[
A=\limsup_{R\to\infty}\frac1{2R}
\int_{[-R,R]\times\R}|F|^2<\infty.
\]
Average translates of this single pair:
\[
\mathsf Q_R=\frac1{2R}\int_{-R}^R
\delta_{(\theta_x\gamma,\theta_xF)}\dd x.
\]
No field is selected as a measurable function of the configuration.

For the field coordinate use $H^{-1}_{\mathrm{loc}}(\R^2;\R^2)$:
a distribution belongs to this space if its product with every smooth
compactly supported cutoff belongs to $H^{-1}(\R^2;\R^2)$, the dual
of $H^1(\R^2;\R^2)$. The local topology is generated by these cutoff
norms on an exhausting sequence of bounded domains. For every $a,b>0$, Fubini gives
\[
\E_{\mathsf Q_R}
\int_{[-a,a]\times[-b,b]}|F|^2
\le \frac{2a}{2R}
\int_{[-R-a,R+a]\times\R}|F|^2,
\qquad
\limsup_{R\to\infty}\E_{\mathsf Q_R}
\int_{[-a,a]\times[-b,b]}|F|^2\le 2aA.
\]
Markov's inequality on an exhausting sequence of boxes and the compact
embedding of bounded local $L^2$ sets into $H^{-1}_{\mathrm{loc}}$
give tightness of the field laws. The configuration marginals converge
to $P$, hence are tight as well. A subsequence therefore has a joint
weak limit $\mathsf Q$, whose configuration marginal is exactly $P$.
The local $L^2$ norm, extended to $+\infty$ outside $L^2$, is lower
semicontinuous for distributional convergence. The same bounds show
that $\mathsf Q$ is supported on locally square-integrable fields.

For each fixed horizontal shift $h$, translating the averaging interval
changes $\mathsf Q_R$ in total variation by at most $|h|/R$.
Translations are continuous in both coordinate topologies, so
$\mathsf Q$ is stationary. For any $u\in C_c^\infty(\R^2)$, the maps
\[
F\longmapsto\langle F,\nabla u\rangle,
\qquad
\gamma\longmapsto
\int\!\left(\int u(p+z_1,z_2)\rho_\eta(\dd z)\right)\gamma(\dd p)
\]
are continuous in these topologies. The inner integral in the second
map is a continuous compactly supported function of $p$. Thus
\eqref{eq:regularized-div} passes to the joint limit, first on a
countable dense family of tests and then on all tests.
Finally the analogous averaging bound for $(0,1)\times(-b,b)$ has
limsup at most $A$. Lower semicontinuity followed by monotone
convergence as $b\to\infty$ proves $\calE\le A$ in
\eqref{eq:field-energy}.
\end{proof}

In the rest of this subsection, an expectation involving $F$ is under
the fixed joint law $\mathsf Q$; for configuration observables it is
under its marginal $P$. We isolate integration by parts with a
configuration-dependent test. For a stationary process of intensity one
with finite local second moments, let $\sigma=\gamma-\Leb$ be the
centered random measure on $\R$. Its covariance spectral measure,
denoted by $\Spec$ (not the relative entropy $S(\mu\mid\omega)$), is
normalized by
\begin{equation}\label{eq:covariance-spectrum}
\E_P|\langle\sigma,f\rangle|^2
=\int_\R|\widehat f(k)|^2S(\dd k),\qquad
\widehat f(k)=\int_\R e^{-ikx}f(x)\dd x.
\end{equation}
The identity is first understood for smooth compactly supported test
functions $f$. Here $\Spec$ is a symmetric positive Radon measure and
need not have a density. Finite local second moments make $\sigma$ a stationary random
tempered distribution, so this definition also applies to Schwartz
tests by approximation.

\begin{samepage}
\begin{lemma}[Stationary random-test integration by parts]
\label{lem:stationary-ibp}
Suppose $P$ has intensity one and finite local second moments, and
$(\gamma,F)$ has a stationary joint law $\mathsf Q$ satisfying
\eqref{eq:regularized-div} and \eqref{eq:field-energy} at fixed $\eta$.
For real even $\chi\in C_c^\infty(\R\setminus\{0\})$, define
\[
h_y(x)=\frac1{2\pi}\int_\R
e^{ikx}\chi(k)e^{-|k||y|}\dd k,
\qquad \varphi(x,y)=\langle\sigma,h_y(x-\cdot)\rangle.
\]
Then
\begin{equation}\label{eq:test-energy}
\E\int_{[0,1]\times\R}|\nabla\varphi|^2
=2\int_\R |k|\chi(k)^2S(\dd k),
\end{equation}
and
\begin{equation}\label{eq:test-pairing}
\E\int_{[0,1]\times\R}F\cdot\nabla\varphi
=2\pi\int_\R\chi(k)a_\eta(k)\Spec(\dd k),
\end{equation}
where
\[
a_\eta(k)=\frac1{2\pi}\int_0^{2\pi}
e^{ik\eta\cos\theta-|k|\eta|\sin\theta|}\dd\theta.
\]
\end{lemma}
\end{samepage}

\begin{proof}
Fix $0<c<C<\infty$ such that $\operatorname{supp}\chi\subset
\{c\le|k|\le C\}$. The kernels $h_y$ and their horizontal
derivatives are Schwartz functions, with each Schwartz seminorm bounded
by a constant times $e^{-c|y|/2}$; polynomial factors arising from
spectral differentiation are absorbed in this exponential bound.
The same holds for their first vertical derivatives on each side of
$y=0$. Local point counts grow at most polynomially almost surely,
by stationarity and the first moment bound. Thus the convolutions
define pathwise smooth functions away from $y=0$, continuous across
$y=0$, with a weak first vertical derivative. Their mean is zero.
Applying \eqref{eq:covariance-spectrum} at each $y$ and integrating
$2k^2e^{-2|k||y|}$ over $\R$ gives
\eqref{eq:test-energy}.

Here are explicit cutoffs for \eqref{eq:test-pairing}. Let
$q_\epsilon$ be a nonnegative even smooth mollifier in the vertical
variable, supported in $[-\epsilon,\epsilon]$, with mass one, and set
$\varphi_\epsilon=q_\epsilon*_y\varphi$, $0<\epsilon<1$.
It is smooth in both variables. The compact spectral support gives,
with constants depending on $\chi$ and $\Spec$ on its support,
\begin{align}
\E\bigl(|\varphi_\epsilon(0,y)|^2+
|\nabla\varphi_\epsilon(0,y)|^2\bigr)
&\le C_\chi e^{-c|y|},\label{eq:test-decay}\\
\E\int_\R|\nabla\varphi_\epsilon(0,y)-
\nabla\varphi(0,y)|^2\dd y
&\le C_\chi\epsilon.\label{eq:test-smoothing}
\end{align}
For the latter estimate, outside $|y|\le\epsilon$ the multipliers and
their first derivatives change by $O(\epsilon)e^{-c|y|}$; inside that
strip the first derivatives are uniformly bounded and the strip has
length $2\epsilon$. This accounts for the cusp without assuming
pointwise values of $F$ on the line.

Take $w\in C_c^\infty(\R)$ nonnegative with $\int w=1$, and put
$w_L(x)=L^{-1}w(x/L)$. Take $\vartheta\in C_c^\infty(\R)$ equal to one
on $[-1,1]$ and supported in $[-2,2]$, and put
$\vartheta_H(y)=\vartheta(y/H)$, where $H>\max(1,\eta)$.
The distribution identity can be tested pathwise against
$w_L\vartheta_H\varphi_\epsilon$: it holds simultaneously for all smooth
compactly supported tests on a single event of probability one.
Consequently the identity also holds for this test after the configuration
has been fixed.
Every expected pairing is integrable. On the field side this follows
from \eqref{eq:field-energy} and \eqref{eq:test-decay}; on the charge
side it follows from local second moments and the Schwartz bounds.
In particular no third moment of the counts is used.

Expand the gradient of this product and take expectation. Stationarity
makes the main term independent of $L$, because $\int w_L=1$.
Cauchy--Schwarz and $\int|w_L'|=L^{-1}\int|w'|$ bound the horizontal
cutoff error by
\[
\left|\E\int_{\R^2}w_L'\vartheta_H F_x\varphi_\epsilon\right|
\le \frac{C_\chi\sqrt{\calE}}{L}.
\]
The vertical cutoff error and the omitted tail of the main term are
bounded, respectively, by
\[
\left|\E\int_{\R^2}w_L\vartheta_H' F_y\varphi_\epsilon\right|
\le\frac{C_\chi\sqrt{\calE}}{H}e^{-cH/2},
\qquad
C_\chi\sqrt{\calE}\,e^{-cH/2}.
\]
These constants are uniform in $L,H$ and $0<\epsilon<1$.

For clarity, the expected charge pairing can be computed before
removing either spatial cutoff. Put
\[
b_\epsilon(k,y)=
\int q_\epsilon(s)e^{-|k||y-s|}\dd s.
\]
The covariance identity and stationarity give, as an identity of
measures in $p$,
\[
\E\bigl[\gamma(\dd p)\varphi_\epsilon(p+a,b)\bigr]
=\left(\int e^{ika}\chi(k)b_\epsilon(k,b)\Spec(\dd k)\right)\dd p.
\]
It follows first by testing against a smooth compactly supported
function of $p$ and then by the Schwartz approximation above. Since
$\vartheta_H=1$ on the support of $\rho_\eta$ and $\int w_L(p+a)\dd p=1$,
the right side of the tested divergence identity is exactly
\[
2\pi\int\chi(k)a_{\eta,\epsilon}(k)\Spec(\dd k),
\qquad
a_{\eta,\epsilon}(k)=
\int e^{ikz_1}b_\epsilon(k,z_2)\rho_\eta(\dd z).
\]
The original line background has zero expected pairing since
$\E\varphi_\epsilon=0$. Thus there is no extra contribution from
the difference between the original and smeared backgrounds.

Let $L\to\infty$, then $H\to\infty$, and finally
$\epsilon\downarrow0$, in that order. The first two limits follow
from the displayed cutoff bounds. By \eqref{eq:test-smoothing} the
field pairing changes in the last limit by at most
$C_\chi\sqrt{\calE\epsilon}$. On $\operatorname{supp}\chi$,
$|a_{\eta,\epsilon}-a_\eta|\le C\epsilon$; the charge pairing
therefore converges as well, since $\Spec$ is finite on that compact set.
This proves \eqref{eq:test-pairing}.
\end{proof}

\begin{samepage}
\begin{proposition}[Spectral energy and hyperuniformity]
\label{prop:spectral}
Suppose $P$ is stationary of intensity one, satisfies
\eqref{eq:linear}, and admits a stationary joint law with an electric field as in
Lemma~\ref{lem:lift}. Then
\begin{equation}\label{eq:spectral-variance-conclusions}
\lim_{t\to\infty}\frac{v_P(t)}t=0,
\qquad \int_1^\infty\frac{v_P(t)}{t^2}\dd t<\infty.
\end{equation}
Thus both \eqref{eq:hyper} and \eqref{eq:integrated} hold.
\end{proposition}
\end{samepage}

\begin{proof}
The linear variance bound gives finite local second moments, so define
$\Spec$ by \eqref{eq:covariance-spectrum}. Approximation of interval
indicators by smooth tests gives
\begin{equation}\label{eq:count-spectrum}
v_P(t)=\int_\R\frac{4\sin^2(kt/2)}{k^2}\Spec(\dd k),
\end{equation}
with the continuous value $t^2$ at zero. In particular,
$t^2S(\{0\})\le v_P(t)\le C_Pt$ for $t\ge1$, hence
$\Spec(\{0\})=0$.

Apply Lemma~\ref{lem:stationary-ibp} and Cauchy--Schwarz to obtain
\begin{equation}\label{eq:spectral-duality}
\left|\int_\R\chi(k)a_\eta(k)\Spec(\dd k)\right|^2
\le\frac{\calE}{2\pi^2}
      \int_\R |k|\chi(k)^2\Spec(\dd k).
\end{equation}
The coefficient is obtained by dividing the product of
\eqref{eq:field-energy} and \eqref{eq:test-energy} by $(2\pi)^2$. The function $a_\eta$ is real and
even, is smooth away from zero, and is continuous with $a_\eta(0)=1$.
Choose $\delta\in(0,1)$ so that $a_\eta(k)\ge1/2$ whenever
$|k|<\delta$.
For a smooth even cutoff $0\le\psi\le1$ supported in
$\epsilon<|k|<\delta$, the test
$\chi(k)=a_\eta(k)\psi(k)/|k|$ belongs to
$C_c^\infty(\R\setminus\{0\})$: every integral is finite on
this compact annulus before any limit is taken. Write
\[
A_\psi=\int\frac{a_\eta(k)^2}{|k|}\psi(k)\Spec(\dd k).
\]
The rightmost integral in \eqref{eq:spectral-duality} is at most
$A_\psi$, because $\psi^2\le\psi$. Therefore
\[
A_\psi^2\le\frac{\calE}{2\pi^2}A_\psi,
\qquad A_\psi\le\frac{\calE}{2\pi^2}.
\]
Exhaust the punctured interval by such cutoffs and use monotone
convergence and $a_\eta^2\ge1/4$ to deduce
\begin{equation}\label{eq:low-frequency}
\int_{0<|k|<\delta}\frac{\Spec(\dd k)}{|k|}
\le\frac{2\calE}{\pi^2}<\infty.
\end{equation}
The annular tests establish this integrability; it was not assumed
in constructing the potential or applying integration by parts.

For high frequencies, integrate \eqref{eq:count-spectrum} over
$1\le t\le2$. The integral of $v_P(t)$ is finite. The continuous
function $\int_1^2\sin^2(kt/2)\dd t$ is strictly positive on every
compact annulus and tends to $1/2$ as $|k|\to\infty$; it therefore
has a positive lower bound on $|k|\ge\delta$. Tonelli yields
\begin{equation}\label{eq:high-frequency}
\int_{|k|\ge\delta}k^{-2}\Spec(\dd k)<\infty.
\end{equation}

Another application of Tonelli, now using
\[
\frac4{k^2}\int_1^\infty\frac{\sin^2(kt/2)}{t^2}\dd t
\le
\begin{cases}
C/|k|,&0<|k|<\delta,\\
4/k^2,&|k|\ge\delta,
\end{cases}
\]
proves \eqref{eq:integrated}. Finally, for $t\ge1$ the integrand
$4\sin^2(kt/2)/(tk^2)$ of $v_P(t)/t$ tends to zero at every
$k\ne0$. On $0<|k|<\delta$ it is bounded by $C/|k|$, using
$\sin^2 u\le\min(u^2,1)$; on $|k|\ge\delta$ it is bounded by
$4/k^2$. Equations \eqref{eq:low-frequency} and
\eqref{eq:high-frequency}, and the absence of an atom at zero,
permit dominated convergence. This proves \eqref{eq:hyper}.
\end{proof}

\subsection{The finite-energy implication}

\begin{samepage}
\begin{proposition}[Variance hypotheses for finite-energy components]
\label{prop:energy-discrepancy}
Let $P$ be an ergodic stationary canonical logarithmic DLR solution
with $\E_PW<\infty$, where $W$ is given by
\eqref{eq:renormalized-energy}. Then $P$ has intensity one, and there is $C_P<\infty$ such that
\begin{equation}\label{eq:finite-energy-variance-conclusions}
v_P(t)\le C_Pt\quad(t>0),\qquad
\lim_{t\to\infty}\frac{v_P(t)}t=0,\qquad
\int_1^\infty\frac{v_P(t)}{t^2}\dd t<\infty.
\end{equation}
These are exactly \eqref{eq:linear}--\eqref{eq:integrated}; $C_P$ may
depend on $P$.
\end{proposition}
\end{samepage}

\begin{proof}
Lemma~\ref{lem:energy-input} gives intensity one and the linear
variance bound. Because $W$ is bounded below, finite expectation
implies $W<\infty$ almost surely. Lemma~\ref{lem:lift} constructs the
stationary joint distribution of the configuration and field; Proposition~\ref{prop:spectral} then proves
the sublinear and integrated variance bounds. In particular, the
integrated bound is a conclusion of the spectral argument, not an
imported consequence of the linear discrepancy estimate alone.
\end{proof}

\begin{proof}[Proof of Theorem~\ref{thm:finite-energy}]
By the canonical DLR ergodic decomposition in
\cite[Proposition~3.13]{DHLM}, there are a probability space
$(\mathcal A,\pi)$ and ergodic stationary canonical DLR solutions
$P_\alpha$ such that
\begin{equation}\label{eq:ergodic-decomposition}
P(f)=\int_{\mathcal A}P_\alpha(f)\pi(\dd\alpha)
\qquad\text{for every bounded measurable }f.
\end{equation}
Here the DLR specification is the conditional-density specification
described in Section~1; its preservation under this decomposition
uses the probabilistic argument of that proposition, independently
of any numerical convention for the energy.
With $C_0$ as in Lemma~\ref{lem:energy-input}, Tonelli's theorem gives
\[
\int_{\mathcal A}\E_{P_\alpha}(W+C_0)\pi(\dd\alpha)
 =\E_P(W+C_0)<\infty.
\]
Hence $\E_{P_\alpha}W<\infty$ for $\pi$-almost every $\alpha$.
Proposition~\ref{prop:energy-discrepancy} puts each such component in
the class of Theorem~\ref{thm:comparison}. Hence all these components
coincide with a single law.

The two existence inputs are used separately. The canonical DLR
property of $\Sine_\beta$ is \cite[Theorem~2.1]{DHLM}.
Its finite expected energy follows from \cite[Corollary~1.2]{LS}:
$\Sine_\beta$ minimizes a free energy with finite minimum, and the
specific relative entropy in that functional is nonnegative.
Equation~\eqref{eq:energy-bridge} converts this statement to
$\E_{\Sine_\beta}W<\infty$ in our convention.
Apply the same ergodic decomposition to $\Sine_\beta$. All its
components belong to the singleton class of
Theorem~\ref{thm:comparison}, so their common law is $\Sine_\beta$.
Every candidate $P$, being a mixture of this law, equals
$\Sine_\beta$.
\end{proof}


\section{Relation to other uniqueness arguments}

The finite-energy scope of Theorem~\ref{thm:finite-energy} agrees with
\cite[Theorem~1.1]{Assiotis}, by \eqref{eq:assiotis-energy} and the
DLR comparison below. Both results allow nonergodic stationary laws.
The two proofs share a passage from entropy to averaged gap statistics
and then to Palm identification. More specifically,
\cite[Proposition~7.3]{Assiotis} obtains an exterior-averaged relative
entropy of order $o(L)$ against the matched circular $\beta$ ensemble.
Cyclic-gap transport then distributes that cost among the gaps
\cite[Lemma~8.3]{Assiotis}, and the circular ensemble's local and Palm
limits identify the process \cite[Sections~8.3--8.6]{Assiotis}.
The present proof uses the static Brascamp--Lieb and logarithmic Sobolev
estimate of Proposition~\ref{prop:gap-entropy}, together with force
tightness and weak quadratic confinement. It compares both candidates
with one lattice-boundary law and cancels the common reference.
The linear and sublinear variance bounds are the discrepancy inputs
recorded in \cite[Proposition~2.7]{Assiotis}; the integrated variance
condition \eqref{eq:integrated} is the additional explicit sufficient
input used by our force-envelope argument and proved here from energy.

\begin{remark}[Agreement of the DLR conventions]\label{rem:dlr-conventions}
On the almost-sure regularity event of Lemma~\ref{lem:as-regularity},
the interval formulation used here extends to the bounded Borel
resampling sets in \cite[Definitions~2.1--2.2]{Assiotis}.
To check its uniformity requirement, fix a bounded Borel set
$\Lambda\subset[-B,B]$ with $B>0$ and $|\Lambda|>0$, put
$n=N_\Lambda(\gamma)$, and define
\[
\Conf_n(\Lambda)=\{\eta\in\Conf(\R):
                 \eta(\Lambda)=n,\ \eta(\Lambda^c)=0\}.
\]
For $\eta\in\Conf_n(\Lambda)$, write
\[
M_{\Lambda,R}(\eta,\gamma)
=-\sum_{\substack{p\in\gamma\setminus\Lambda\\|p|\le R}}
\int_\Lambda\log|x-p|\,(\eta-\gamma_\Lambda)(\dd x)
\]
for the relative exterior interaction difference at radius $R$.
For $S>R>2B$, Taylor expansion and cancellation of the total mass give,
uniformly over all $n$-point configurations $\eta$ in $\Lambda$,
\[
|M_{\Lambda,S}(\eta,\gamma)-M_{\Lambda,R}(\eta,\gamma)|
\le 2nB\left|\sum_{\substack{p\in\gamma\\R<|p|\le S}}\frac1p\right|
+CnB^2\sum_{\substack{p\in\gamma\\|p|>R}}\frac1{p^2}.
\]
Density one makes the second series summable, and the reciprocal
principal value makes the first term Cauchy. Thus the relative interaction differences converge
uniformly on the set of all $n$-point configurations in $\Lambda$, including along the radii $m/2$, $m\in\N$, used
in \cite{Assiotis}.

Choose a bounded interval containing $\Lambda$ and condition its
canonical Gibbs law further on the points outside $\Lambda$ and on
$N_\Lambda$. Disintegration gives precisely the kernel of
\cite[Definition~2.1]{Assiotis}. Its normalizer is positive and finite:
the distant potential is bounded after fixing its additive constant,
while the finitely many nearby logarithmic interaction factors are
bounded above and positive Lebesgue-almost everywhere on the bounded
domain. The binomial reference measure is the law of $n$ independent points
uniformly distributed in $\Lambda$. Its ordered-coordinate density is
$n!/\Leb(\Lambda)^n$, a constant which cancels in normalization. Subtracting the interaction difference to a
fixed reference configuration with the same exterior and particle number removes dependence on the
original interior configuration and gives exterior-and-count
measurability. Limits of the finite-cutoff kernels, or a conditional-law
version with a countable determining class, give a measurable version of the conditional kernel, up to null sets,
as required by the formulation of \cite[Definition~2.2]{Assiotis}. For particle number zero, the kernel is the point mass on the empty
configuration. Conversely, that definition immediately includes
interval resampling. Assiotis initially allows multiplicities, but
\cite[Lemma~2.8]{Assiotis} makes every canonical DLR law simple.

For the finite-energy class, the required regularity holds on almost
every ergodic component by Proposition~\ref{prop:energy-discrepancy}
and Lemma~\ref{lem:as-regularity}, hence also under their mixture.
This verifies agreement of the DLR classes used in the two
finite-energy statements.
\end{remark}

At $\beta=2$, the deterministic theorem of \cite[Theorem~1.3]{KMD}
provides another route. The reciprocal principal value and density-one
conditions needed there follow from the discrepancy bounds, so its
kernel convergence applies to every admissible candidate boundary.
DLR and convergence of compactly supported Laplace functionals then
identify the process.

For general $\beta$, the main theorem of the arXiv version of
\cite{BEY} specified in the bibliography concerns a fixed real analytic uniformly convex potential
on the line, and its local comparison is combined with strong
particle-location estimates. A direct application to arbitrary
random DLR exterior potentials would require checking assumptions not
provided by macroscopic potential convergence alone.
The present proof uses the static convexity inequalities from that
framework and compares with one common reference law. In particular,
it establishes neither a separate universality theorem for all
deterministic boundaries nor a limiting theorem for
$\omega_{K,T_K}$.

The distinction between the two improved discrepancy inputs is useful:
$v_P(t)=o(t)$ makes the average quadratic displacement $o(K)$ per
particle, whereas $\int_1^\infty v_P(t)t^{-2}\dd t<\infty$ supplies
the half-line force envelopes. The linear bound alone gives neither
of these two improvements. The weak confinement can be chosen between
the quadratic displacement scale and the block length precisely
because the improvement is little-$o$; no fixed power rate is needed.

\appendix
\section{Convex inequalities on the singular simplex}
\label{app:convex}

We record the approximation used in Proposition~\ref{prop:gap-entropy}.
Let $\nu$ be a probability measure with a positive continuous Lebesgue
density inside $\Omega_K$. Define its weighted weak Sobolev space by
\begin{equation}\label{eq:weighted-sobolev}
H^1(\nu)=\left\{f\in H^1_{\rm loc}(\Omega_K):
\int_{\Omega_K}(f^2+|\nabla f|^2)\dd\nu<\infty\right\}.
\end{equation}
Here $\nabla f$ is the distributional weak gradient.
This definition imposes no boundary trace condition.
For a nonnegative function $h$, write
\[
\Ent_\nu(h)=\int h\log h\dd\nu
 -\left(\int h\dd\nu\right)\log\left(\int h\dd\nu\right),
\]
with the usual extended-value convention.

\begin{samepage}
\begin{lemma}[Removal of the gap truncation]\label{lem:convex-approx}
Suppose $V\in C^2(\Omega_K)$, $0<Z=\int_{\Omega_K}e^{-V}<\infty$,
and $V''\ge\rho\Id_K$ for some $\rho>0$.
Put $\nu=Z^{-1}e^{-V}\dd u$. For every $f\in H^1(\nu)$,
\begin{align}
\Var_\nu f
&\le\int\nabla f^\top(V'')^{-1}\nabla f\dd\nu,
\label{eq:appendix-bl}\\
\Ent_\nu(f^2)
&\le\frac2\rho\int|\nabla f|^2\dd\nu
\qquad(f\in H^1(\nu)).\label{eq:appendix-lsi}
\end{align}
In particular, for $\mu=h\nu$ with $\sqrt h\in H^1(\nu)$,
$S(\mu\mid\nu)\le I(\mu\mid\nu)/(2\rho)$, where
$I=4\int|\nabla\sqrt h|^2\dd\nu$.
These statements apply to $V=\Phi_{K,T}-sF_K$ throughout the tilt
range in Section~\ref{sec:entropy}, with the Hessian bounds stated
there and constants unaffected by removal of the truncation.
\end{lemma}
\end{samepage}

\begin{proof}
For $0<\delta<1$ define the nonempty convex simplex
\[
\Omega_{K,\delta}=\{u:g_i(u)>\delta\text{ for }0\le i\le K\},
\qquad
\nu_\delta=\frac{\1_{\Omega_{K,\delta}}\nu}
 {\nu(\Omega_{K,\delta})}.
\]
Its closure lies strictly inside $\Omega_K$, so the density there
is bounded above and below by positive constants and $V$ is $C^2$
on a neighborhood of the closure. The bounds may depend on $\delta$; the convexity
constant $\rho$ does not. When applying a smooth-domain argument below,
one may first mollify $V$ on this neighborhood. Mollification preserves
$V''\ge\rho\Id_K$, and convergence in $C^2$ on the closure allows
passage back to the stated $C^2$ potential.

We specify the convex-domain inequalities being used. For a smooth
bounded convex domain $D$ and $\nu_D\propto e^{-V}\1_D$, the
Brascamp--Lieb and Bakry--\'Emery inequalities have precisely the
forms \eqref{eq:appendix-bl}--\eqref{eq:appendix-lsi}, with $D$
as the integration domain; see \cite{BL,BE} for the underlying
inequalities. Their boundary extension can be seen directly with
the Neumann generator $\mathcal L=\Delta-\nabla V\cdot\nabla$.
Let $n$ denote the outward unit normal on $\partial D$.
If $-\mathcal Lw=f-\nu_D(f)$ and $\partial_nw=0$, the integrated
Bochner identity and convexity of $\partial D$ give
\[
\int(\mathcal Lw)^2\dd\nu_D
\ge\int\nabla w^\top V''\nabla w\dd\nu_D.
\]
The omitted nonnegative terms are
$\int\|D^2w\|_{\rm HS}^2\dd\nu_D$ and the boundary integral of the
second fundamental form evaluated on the tangential gradient. Here
$\|D^2w\|_{\rm HS}^2=\sum_{i,j}|\partial_{ij}w|^2$ is the squared
Hilbert--Schmidt norm of the Hessian. Integration by parts
and weighted Cauchy--Schwarz then give
\[
\Var_{\nu_D}f=\int\nabla f\cdot\nabla w\dd\nu_D
\le\left(\int\nabla f^\top(V'')^{-1}\nabla f\dd\nu_D\right)^{1/2}
   (\Var_{\nu_D}f)^{1/2},
\]
which proves the variance statement. For the Neumann semigroup
$P_t$, entropy dissipation for positive $h$ and the same Bochner
identity yield
\[
-\frac{\dd}{\dd t}\Ent_{\nu_D}(P_th)
 =\mathscr I_D(P_th),\qquad
\frac{\dd}{\dd t}\mathscr I_D(P_th)\le-2\rho \mathscr I_D(P_th),\qquad
\mathscr I_D(h)=\int\frac{|\nabla h|^2}{h}\dd\nu_D.
\]
Integrating to equilibrium gives
$\Ent_{\nu_D}(h)\le \mathscr I_D(h)/(2\rho)$.
Apply this to $h=f^2+\varepsilon$ and let $\varepsilon\downarrow0$
to obtain the stated factor $2/\rho$. Smooth approximation on a
smooth bounded domain extends both inequalities to its ordinary
$H^1$ space. This step uses only weights bounded above and below,
and makes no assertion about density at collision faces.

Corners of $\Omega_{K,\delta}$ can be removed without changing
the potential or its Hessian: the sublevel domains
\[
D_m=\left\{u\in\Omega_{K,\delta}:
 \sum_{i=0}^K\frac1{g_i(u)-\delta}<m\right\}
\]
are smooth bounded convex domains for $m$ above the minimum of
the displayed strictly convex function, and exhaust
$\Omega_{K,\delta}$. Applying the preceding inequalities on
$D_m$ and passing to the limit gives them for $\nu_\delta$.
Every fixed $f\in H^1(\nu)$ restricts to ordinary $H^1$ on these
domains and on $\Omega_{K,\delta}$, by local comparability of the
weights. Variance passes by convergence of its two moments, and
the right sides pass by monotone convergence of the unnormalized
integrals. The entropy passage is justified by the exact formula
below, used first with $D_m$ and then with $\Omega_{K,\delta}$.

Indeed, set $z_\delta=\nu(\Omega_{K,\delta})$ and
$a_\delta=\int_{\Omega_{K,\delta}}f^2\dd\nu$. Then
\[
z_\delta\Ent_{\nu_\delta}(f^2)
=\int_{\Omega_{K,\delta}}f^2\log f^2\dd\nu
 -a_\delta\log(a_\delta/z_\delta).
\]
As $\delta\downarrow0$, $z_\delta\to1$ and
$a_\delta\to\nu(f^2)$. The negative part of $f^2\log f^2$ is
bounded by $1/e$, while its positive part passes by monotone
convergence. Thus the left side converges to
$\Ent_\nu(f^2)$, even if initially that quantity is infinite.
Multiplying the truncated logarithmic Sobolev inequality by
$z_\delta$ and taking the limit proves
\eqref{eq:appendix-lsi} for the full weak Sobolev domain directly.
The same restriction argument proves \eqref{eq:appendix-bl}.

For $V=\Phi_{K,T}-sF_K$, the normalizing integral is finite because
$F_K$ is bounded. Throughout $|s|\le c_GK$ its Hessian satisfies
both $V''\ge\frac12\Phi_{K,T}''$ and $V''\ge T^{-1}\Id_K$ on every
truncated domain. In particular, the pointwise bound
$\nabla F_K^\top(V'')^{-1}\nabla F_K\le2A_1/K$ survives the
limit unchanged. Taking $s=0$, $\rho=T^{-1}$ and $f=\sqrt h$
proves \eqref{eq:lsi} in exactly the Fisher-information convention
used in Section~\ref{sec:entropy}. All integrations by parts took
place at strictly positive gaps. No inverse-gap integrability or
boundary trace at a collision is required, for any $\beta>0$.
\end{proof}

\begin{thebibliography}{DHLM19}

\bibitem[Ass26]{Assiotis}
Theodoros Assiotis,
\emph{Uniqueness for DLR equations of $\mathsf{Sine}_\beta$}.
Preprint arXiv:2609.06832v1, 6 September 2026.
The theorem, definition, and section numbers cited here refer to this version.
\url{https://arxiv.org/abs/2609.06832v1}.

\bibitem[BL76]{BL}
Herm Jan Brascamp and Elliott H.~Lieb,
\emph{On extensions of the Brunn--Minkowski and Pr\'ekopa--Leindler
theorems, including inequalities for log concave functions, and with
an application to the diffusion equation}.
\emph{Journal of Functional Analysis} \textbf{22} (1976), no.~4,
366--389. \href{https://doi.org/10.1016/0022-1236(76)90004-5}
{doi:10.1016/0022-1236(76)90004-5}.

\bibitem[BE85]{BE}
Dominique Bakry and Michel \'Emery,
\emph{Diffusions hypercontractives}.
\emph{S\'eminaire de probabilit\'es XIX}, Lecture Notes in Mathematics
\textbf{1123}, Springer, 1985, 177--206.
\href{https://www.numdam.org/item/SPS_1985__19__177_0/}
{Numdam full text}.


\bibitem[BEY12]{BEY}
Paul Bourgade, L\'aszl\'o Erd\H{o}s and Horng-Tzer Yau,
\emph{Universality of General $\beta$-Ensembles}.
\emph{Duke Mathematical Journal} \textbf{163} (2014), no.~6,
1127--1190. The theorem numbering used here refers to
arXiv:1104.2272v5, 5 February 2012.
\url{https://arxiv.org/abs/1104.2272v5}.

\bibitem[DHLM19]{DHLM}
David Dereudre, Adrien Hardy, Thomas Lebl\'e and Myl\`ene Ma\"ida,
\emph{DLR equations and rigidity for the Sine-beta process}.
Preprint arXiv:1809.03989v2, 15 November 2019
(title-page date 18 November 2019).
Theorem numbering and the displayed convention discussed in
Remark~\ref{rem:cutoff} refer to this version.
\url{https://arxiv.org/abs/1809.03989v2}.

\bibitem[KMD19]{KMD}
Arno B.~J. Kuijlaars and Erwin Mi\~na-D\'iaz,
\emph{Universality for conditional measures of the sine point process}.
\emph{Journal of Approximation Theory} \textbf{243} (2019), 1--24.
References here use arXiv:1703.02349v2, 21 March 2019.
\url{https://arxiv.org/abs/1703.02349v2}.

\bibitem[LS17]{LS}
Thomas Lebl\'e and Sylvia Serfaty,
\emph{Large deviation principle for empirical fields of Log and Riesz gases}.
\emph{Inventiones Mathematicae} \textbf{210} (2017), no.~3, 645--757.
\href{https://doi.org/10.1007/s00222-017-0738-0}
{doi:10.1007/s00222-017-0738-0}.
Section and equation conventions used here refer to the corrected
manuscript dated 9 May 2017,
\url{https://math.nyu.edu/~serfaty/LDPcorrection.pdf}.

\end{thebibliography}
\end{document}
